\documentclass[11pt]{article}

\usepackage[T1]{fontenc}
\usepackage[margin=1in]{geometry}
\usepackage{lmodern}
\usepackage{microtype}
\usepackage{amsmath,amssymb,mathtools}
\usepackage{amsthm}
\usepackage{aliascnt}
\usepackage{booktabs}
\usepackage{tabularx}
\usepackage[font=small,labelfont=bf]{caption}
\usepackage[section]{placeins}
\usepackage{float}
\usepackage{array}
\usepackage{xcolor}
\usepackage{enumitem}
\usepackage{needspace}
\usepackage[hidelinks]{hyperref}
\usepackage[nameinlink,noabbrev]{cleveref}
\usepackage{url}
\usepackage[backend=biber,style=numeric,sorting=nyt,maxbibnames=99]{biblatex}
\addbibresource{references.bib}


\newtheorem{theorem}{Theorem}

\newaliascnt{lemma}{theorem}

\newtheorem{lemma}[lemma]{Lemma}
\aliascntresetthe{lemma}

\newaliascnt{proposition}{theorem}

\newtheorem{proposition}[proposition]{Proposition}
\aliascntresetthe{proposition}

\newaliascnt{corollary}{theorem}

\newtheorem{corollary}[corollary]{Corollary}
\aliascntresetthe{corollary}
\crefname{theorem}{Theorem}{Theorems}
\crefname{lemma}{Lemma}{Lemmas}
\crefname{proposition}{Proposition}{Propositions}
\crefname{corollary}{Corollary}{Corollaries}

\definecolor{deepblue}{RGB}{23,69,124}
\definecolor{lightblue}{RGB}{242,247,252}
\definecolor{darkred}{RGB}{145,32,32}
\ifdefined\ANONYMOUS

\newcommand{\pdfpaperauthor}{Anonymous}
\else

\newcommand{\pdfpaperauthor}{Justin Thaler, Kai Zhang, Scott Kominers, Quang Dao}
\fi
\hypersetup{
  colorlinks=true,
  linkcolor=deepblue,
  citecolor=deepblue,
  urlcolor=deepblue,
  pdftitle={Public-Matrix Symmetry Yields Faster Forgery Attacks on Odd-Characteristic SNOVA},
  pdfauthor={\pdfpaperauthor},
  pdfsubject={Cryptanalysis of SNOVA Version 2.3},
  pdfkeywords={SNOVA, multivariate signatures, forgery attack, NIST PQC}
}
\setlist{leftmargin=*,topsep=3pt,itemsep=1.5pt,parsep=0pt,partopsep=0pt}


\newcommand{\F}{\mathbb F}

\newcommand{\Mat}{\operatorname{Mat}}

\newcommand{\Sym}{\operatorname{Sym}}

\newcommand{\rank}{\operatorname{rank}}

\newcommand{\Span}{\operatorname{span}}

\newcommand{\transpose}{\mathsf T}

\newcommand{\SNOVA}{\textsf{SNOVA}}

\newcommand{\bits}{\text{ bits}}

\newcommand{\LoneModelRange}{$2^{128.82}$--$2^{129.44}$}

\newcommand{\LthreeModelRange}{$2^{171.69}$--$2^{183.44}$}

\newcommand{\LfiveModelRange}{$2^{214.12}$--$2^{235.19}$}

\newcommand{\LoneAXNRange}{$2^{132.17}$--$2^{133.75}$}

\newcommand{\LthreeAXNRange}{$2^{175.04}$--$2^{187.74}$}

\newcommand{\LfiveAXNRange}{$2^{217.48}$--$2^{239.50}$}

\newcommand{\LoneNANDRange}{$2^{133.04}$--$2^{134.67}$}

\newcommand{\LthreeNANDRange}{$2^{175.91}$--$2^{188.67}$}

\newcommand{\LfiveNANDRange}{$2^{218.35}$--$2^{240.42}$}

\newcommand{\AXNNominalHeadroomRange}{9.25--54.52}

\newcommand{\AXNExpectedAESGapRange}{8.13--52.86}

\newcommand{\AXNFullAESGapRange}{9.13--53.86}

\newcommand{\NANDNominalHeadroomRange}{8.33--53.65}

\newcommand{\NANDExpectedAESGapRange}{9.03--53.83}

\newcommand{\NANDFullAESGapRange}{10.03--54.83}

\newcommand{\attacktitle}{Public-Matrix Symmetry Yields Faster Forgery
  Attacks on Odd-Characteristic \texorpdfstring{\SNOVA}{SNOVA}}

\title{\attacktitle}
\ifdefined\ANONYMOUS
  \author{Anonymous submission}
\else
  \author{%
    Justin Thaler\thanks{a16z crypto research and Georgetown University} \and
    Kai Zhang\thanks{MIT} \and
    Scott Kominers\thanks{Harvard University and a16z crypto research} \and
    Quang Dao\thanks{Carnegie Mellon University and LayerZero Labs}%
  }
\fi
\date{Draft of 27 July 2026}

\begin{document}
\setlength{\emergencystretch}{2em}
\raggedbottom
\maketitle

\ifdefined\DISCLOSURE
\begin{center}
\fcolorbox{darkred}{white}{\parbox{0.91\textwidth}{\centering\small
\textbf{Confidential coordinated disclosure.} This draft is intended for the
\SNOVA{} submission team and NIST.\@ Please do not redistribute before the
coordinated public-disclosure date.}}
\end{center}
\fi

\begin{abstract}
At its algebraic core, a \SNOVA{} forgery requires solving a public system of
multivariate quadratic (MQ) equations; the resulting assignment must also pass
the verifier's rejection checks.  We show that the public $q=19$ design
represents many of its quadratic features twice.  A feature is labeled by an \emph{ordered} pair of
powers from the public algebra $A=\F_q[S]$, but symmetry of the public base
matrices makes the polynomial unchanged when the two labels are reversed.
Thus, if $d=\dim_{\F_q}A$, the $d^2$ labeled power-pair features of each
underlying public form factor through at most $\binom{d+1}{2}$ unordered
features.  In the public parameter sets this
reduces $16$ labels to $10$ at block size $\ell=4$, and $4$ labels to $3$
at block size $\ell=2$.  The identification occurs before the public outer map and holds for
every symmetric public key; it is not a weak-key event.

For each of the nine public $q=19$ parameter sets in \SNOVA{}
Version~2.3, we combine this symmetry quotient with exact affine elimination to
construct a publicly checkable reduction to a smaller forgery system, without
dropping verifier equations.  For $\ell=4$, the chosen
column relation preserves the four-block structure used by special XL.\@  For
$\ell=2$, fixed-skew and target-filtered relations ensure that recovered roots
also pass the verifier's block-symmetry check.  The quotient, affine reduction,
rejection identities, target-sampling argument, probability bounds, and final
verification wrapper are exact.  Under the stated global-rank conditions and
the cited production-scale MQ models, the modeled expected operation counts at
NIST Levels~I, III, and V are, respectively, \LoneModelRange,
\LthreeModelRange, and \LfiveModelRange.  Replacing the field
arithmetic charged by those models with validated two-input AND/XOR/XNOR
circuits gives dominant-arithmetic gate counts \LoneAXNRange,
\LthreeAXNRange, and \LfiveAXNRange.  These are not complete attack
circuits: the global rank
spectra, production solver behavior, control logic, memory, and data movement
remain the principal qualifications.
\end{abstract}

\section{Introduction}\label{sec:intro}

\paragraph{The cryptanalytic task.}
A multivariate signature scheme publishes a polynomial map of total degree
at most two,
\[
  \mathcal P:\F_q^{N}\longrightarrow\F_q^{M}.
\]
For a message with hash-derived target $t\in\F_q^M$, the algebraic core of a
forgery is to find $x$ satisfying $\mathcal P(x)=t$; any additional format or
rejection checks must also pass.  This is an MQ problem: $M$ equations of
degree at most two in the coordinates of $x$.  Linear and constant terms are
allowed; throughout the paper, ``quadratic'' refers to the maximum degree.
The apparent equation and variable counts are only a first approximation to
security.  Public algebraic structure can make different-looking quadratic
features identical, can expose linear equations, or can split the remaining
variables into blocks that a specialized solver can exploit.

\SNOVA{} is one of the nine digital-signature candidates that NIST advanced to
the third round of its additional post-quantum signature process in
May~2026~\cite{NISTThirdRound2026}.  It stores the signature variables in small
matrices and builds its public quadratics from pairs of powers of a public
matrix $S$.  Beullens showed
that an attacker can impose an affine relation among the signature columns and
thereby reduce forgery to a smaller underdetermined MQ system
\cite{Beullens2024SNOVA}.  We find an additional, deterministic reduction
inside the quadratic features themselves.

\paragraph{The hidden duplication.}
Let $P_i$ be a public symmetric base matrix, let
$A=\F_q[S]$, and write $\Sigma_\xi=I_n\otimes\xi$ for $\xi\in A$.  The feature
with left label $\xi$ and right label $\eta$ is
\[
  q_{i,\xi,\eta}(u)
  =u^\transpose\Sigma_\xi P_i\Sigma_\eta u.
\]
Because the matrices are symmetric and the result is a scalar,
\begin{equation}\label{eq:intro-swap}
  q_{i,\xi,\eta}(u)=q_{i,\eta,\xi}(u).
\end{equation}
The verifier begins from ordered labels, but the polynomial depends only on
the corresponding unordered pair.  We call the operation of identifying
these duplicate labels the \emph{symmetry quotient}.

For a two-dimensional toy algebra with basis $\{1,S\}$, the four labels
\[
 (1,1),\qquad(1,S),\qquad(S,1),\qquad(S,S)
\]
produce only three polynomials because the middle two are equal.  More
generally, if $\dim_{\F_q}A=d$, then at most
\begin{equation}\label{eq:intro-K}
  K=m_1\binom{d+1}{2}
\end{equation}
independent power-pair features remain across the $m_1$ base forms, rather
than $m_1d^2$.  Thus each $\ell=4$ base form drops from $16$ ordered labels to
$10$ unordered labels; each $\ell=2$ base form drops from $4$ to $3$.  The
loss occurs before the public outer linear map---the final public mixing of
these features into verifier outputs---so later randomization cannot recreate
the discarded directions.  It holds for every symmetric public key
and every attacker-chosen column relation.

\begin{center}
\fcolorbox{deepblue}{lightblue}{\parbox{0.94\textwidth}{\small
\textbf{Core idea in plain language.}  The verifier assigns two names to the
same quadratic polynomial whenever it reverses a pair of power labels.  The
attack first removes those duplicate features, then uses ordinary linear
algebra to retain all remaining verifier equations, and finally solves only
the smaller quadratic core.}}
\end{center}

\paragraph{The attack as a data-flow.}
The proof is easiest to follow as a sequence of exact transformations.  The
word \emph{lift} below means the attacker-chosen affine rule that turns one
column vector into a full signature candidate.

\begin{center}
\renewcommand{\arraystretch}{1.25}
\begin{tabularx}{0.96\textwidth}{@{}>{\bfseries\raggedright\arraybackslash}p{0.22\textwidth}>{\centering\arraybackslash}p{0.04\textwidth}X@{}}
Full public verifier & $\longrightarrow$ & impose one affine common-column
relation; quadratic, linear, and constant terms become explicit\\
& $\downarrow$ & \\
Symmetry quotient & $\longrightarrow$ & identify $(\xi,\eta)$ with
$(\eta,\xi)$ and keep only the distinct quadratic features\\
& $\downarrow$ & \\
Exact affine elimination & $\longrightarrow$ & use a left-kernel projection
to expose every output direction unreachable by the quadratics as an ordinary
linear constraint\\
& $\downarrow$ & \\
Structured residual MQ & $\longrightarrow$ & choose a stable four-block slice
for $\ell=4$, or a rejection-compatible relation for $\ell=2$\\
& $\downarrow$ & \\
Candidate signature & $\longrightarrow$ & map back to the original
coordinates and run the unmodified public verifier
\end{tabularx}
\end{center}

Every transformation through the residual MQ system is exact and uses only
public algebra and linear algebra.  The solver model enters only when we
estimate the cost of solving that residual system.  The final public-verifier
check makes the attack Las Vegas: a trial may be discarded, but an invalid
signature is never returned.

\paragraph{Four jobs that a complete attack must perform.}
The equality \eqref{eq:intro-swap} is only the first job.  The construction
also has to preserve all public outputs, ensure that targets have roots, and
make the residual equations compatible with a practical solver:
\begin{enumerate}[label=\textbf{\arabic*.},itemsep=3pt]
\item \textbf{Constrain the signature columns.}  Every signature column is
      written as an affine function of one common vector.  Substitution
      separates the verifier into quotient quadratic features, offset-linear
      terms, and constants.
\item \textbf{Eliminate the affine part without dropping equations.}  Output
      directions outside the quadratic image become linear constraints.  For
      each target satisfying those affine consistency conditions, Gaussian
      elimination leaves a smaller MQ system with exactly the same solutions
      as the constrained verifier.
\item \textbf{Choose a solver-friendly search space.}  For $\ell=4$, a
      structured choice of offsets leaves a large $A$-linear kernel.  Over a
      splitting field this kernel becomes four variable blocks, preserving the
      multidegrees used by special XL.\@  For $\ell=2$, fixed-skew and filtered
      relations handle the verifier's rejection rule directly.
\item \textbf{Retry and verify.}  Exact moment bounds control the
      probabilities of root existence and uniqueness for random target--slice
      pairs.  A candidate recovered by the MQ solver is transformed back and
      checked by the original verifier.
\end{enumerate}

\paragraph{A complete dimension walk-through.}
The Level-I parameter set $(v,o,\ell,r)=(28,5,4,4)$ illustrates where the
reduction comes from: it has 28 vinegar blocks, 5 oil blocks, and blocks with
4 rows and 4 columns.  Here $n=v+o=33$, one stacked signature column has
$N_0=n\ell=132$ coordinates, the verifier has $M=or\ell=80$ scalar outputs,
and there are $m_1=5$ public base-form groups.

\begin{table}[H]
\centering
\caption{The dimensions in the $(28,5,4,4)$ attack.  This table is only a
worked example; the general formulas appear in later sections.}
\label{tab:running-example}
\small
\renewcommand{\arraystretch}{1.12}
\begin{tabularx}{0.96\textwidth}{@{}>{\raggedright\arraybackslash}p{0.28\textwidth}>{\centering\arraybackslash}p{0.18\textwidth}X@{}}
\toprule
Stage & Dimension & Meaning\\
\midrule
Raw public feature vector & $N_{\rm raw}=1{,}280$ & $m_1\ell^2r^2$ labeled slots before
the column relation\\
After the common-column relation & $80$ ordered quadratics & one feature for
each of $m_1\ell^2=5\cdot16$ ordered power pairs\\
After the symmetry quotient & $K=50$ distinct quadratics &
$m_1\binom{\ell+1}{2}=5\cdot10$ unordered power pairs\\
Affine output directions & $M-K=30$ & linear equations that determine an
affine search space\\
Full residual domain & $102$ coordinates & $N_0-(M-K)=132-30$ after affine
elimination\\
Stable kernel & $52$ base-field dimensions & an $A$-linear space of
$A$-dimension $13$ on which offset-linear motion vanishes\\
Selected special-XL slice & $h=40$ & four blocks of $s=10$ variables; hence
$\tau=K-h=10$\\
\bottomrule
\end{tabularx}
\end{table}

The final slice deliberately has fewer variables than equations.  Under full
cross-polar rank, a random target--slice trial then contains a unique
base-field root with probability essentially $19^{-10}$.  Paying roughly
$19^{10}$ retries is worthwhile because reducing the Macaulay matrix from the
full residual domain to four blocks of ten variables saves much more.  This is
why the paper analyzes both the algebraic reduction and the root-retry
probability, rather than quoting an MQ dimension alone.

\paragraph{Headline estimates.}
\Cref{tab:main} reports modeled costs and the differences between them.
``Model'' is the base-2 expected-cost exponent in the \emph{pinned} attack
model, meaning that we fix the cited estimator, its version, and its conventions
exactly as specified in the artifact.  It starts from the cited per-solve MQ
operation formula and, when the solve is repeated, includes
the explicit target--slice retry factor.  Target generation performed before a
single solve is accounted for separately.  ``AXN arith.'' changes only the unit
for the arithmetic charged by that model, replacing its coarse field-operation
factor with validated signed two-input AND/XOR/XNOR (AXN) circuits.  ``Gain over generic'' is the reduction in the Model exponent relative to a pinned generic-MQ
baseline on the same symmetry quotient.  It is not always a solver improvement: for the structured rows it
comes from the stable four-block slice and special XL, whereas for the filtered
Level-V row it comes from the zero-offset filtered lift.  ``Arithmetic budget''
is the nominal NIST reference minus the AXN arithmetic exponent.  It
is room for unpriced overhead in that one metric, not an unconditional
security margin.

\begin{table}[H]
\centering
\caption{Headline conditional estimates for the public Version~2.3 parameter
sets.  The tuple is $(v,o,\ell,r)$; its meaning is given in
\Cref{sec:parameters}.  Model and AXN are base-2 cost exponents; gain over generic
and arithmetic budget are differences of exponents, measured in bits.  V2.3
reproduces the official ``log \#gates'' table.}
\label{tab:main}
\footnotesize
\setlength{\tabcolsep}{2.3pt}
\renewcommand{\arraystretch}{1.08}
\begin{tabular}{@{}c l l r r r r r@{}}
\toprule
Level & Parameters & Mode & V2.3 & Model & \shortstack{AXN\\arith.} &
\shortstack{gain over\\generic} & \shortstack{Arith.\\budget}\\
\midrule
I   & $(28,5,4,4)$  & structured XL & 194 & 128.82 & 132.17 &  7.22 & 10.83\\
I   & $(48,16,2,2)$ & fixed-skew MQ & 164 & 129.44 & 133.75 &  0.00 &  9.25\\
I   & $(28,4,4,5)$  & structured XL & 194 & 128.82 & 132.17 &  7.22 & 10.83\\
\midrule
III & $(40,7,4,4)$  & structured XL & 264 & 171.69 & 175.04 &  8.91 & 31.96\\
III & $(72,24,2,2)$ & fixed-skew MQ & 234 & 183.44 & 187.74 &  0.00 & 19.26\\
III & $(38,5,4,5)$  & structured XL & 239 & 171.69 & 175.04 &  8.91 & 31.96\\
\midrule
V   & $(50,9,4,4)$  & structured XL & 333 & 214.12 & 217.48 & 10.56 & 54.52\\
V   & $(96,32,2,2)$ & filtered MQ   & 303 & 235.19 & 239.50 &  2.65 & 32.50\\
V   & $(52,6,4,6)$  & structured XL & 333 & 214.12 & 217.48 & 10.56 & 54.52\\
\bottomrule
\end{tabular}
\end{table}

For example, the first row's AXN entry $132.17$ means about $2^{132.17}$
charged two-input arithmetic gates under the special-XL model and its
unique-root retry factor.  It does \emph{not} include a complete implementation
of sparse matrix construction, pivot control, memory traffic, or data
movement.

\paragraph{Conditions behind the table.}
The structured-XL rows require two separate production assumptions: every
nonzero output direction must have full cross-polar rank on the chosen slice,
and a planted-solvable special-XL Macaulay matrix (one generated together
with a known root) must have the predicted one-dimensional right kernel.  The fixed-skew Level-I and Level-III rows use
the sufficient all-target root conditions $r_*\ge96$ and $r_*\ge144$.  The
filtered Level-V row uses the accepted-root condition of
\Cref{prop:l2-rejection}, implied there by $r_*\ge240$.  Each $r_*$ condition
is global: it concerns every nonzero output combination, not only the sampled
directions in the artifact.  Lower ranks are handled by explicit retry bounds
in \Cref{sec:slices}.

\begin{center}
\fcolorbox{deepblue}{lightblue}{\parbox{0.94\textwidth}{\small
\textbf{Four levels of claim.}\par\smallskip
\begin{tabularx}{\linewidth}{@{}>{\bfseries}p{0.24\linewidth}X@{}}
Exact mathematics & Symmetry quotient, affine elimination, stable-kernel
construction, fixed-skew identity, rejection tests, uniform-target sampling,
probability inequalities, and final-verifier wrapper.\\[2pt]
Public preflight & Deterministic rank and rejection checks for the chosen key,
relation, offsets, and slice.  A failing choice is discarded before an
expensive MQ solve.\\[2pt]
Finite evidence & Large but nonexhaustive rank campaigns and smaller Macaulay
experiments.  They support, but do not certify, the global production
assumptions.\\[2pt]
Cost model & Published generic-MQ or special-XL arithmetic schedules with an
explicit retry factor and a validated circuit cost for the charged field
primitive.  A complete solver implementation is not priced.
\end{tabularx}}}
\end{center}

The official Version~2.3 values are numerically $30.25$--$115.52$ bits above
the AXN arithmetic column, but the two sets of figures do not come from a
shared complete-circuit methodology.  A cleaner internal comparison recosts a
historical ordered-label affine-column estimator in the same AXN basis; the
resulting per-row reductions are $38.02$--$129.66$ bits
(\Cref{tab:ordered-baseline}).  The unreduced affine-column equations are an
executable system, but \Cref{thm:symmetric-square} proves that their reversed
power labels do not represent independent quadratic features on a symmetric
key.  What is counterfactual is feeding the ordered-label count into a generic
estimator as though those labels described a generic quadratic space.  The
comparison is included only to isolate the numerical effect of recognizing
the quotient and then choosing the selected post-quotient attack mode.

NIST describes its category references across the security metrics it regards
as relevant, rather than as a gate-count-only test~\cite{NISTPQCCriteria}.  We
therefore use $143/207/272$ only as nominal arithmetic reference points.
Establishing an end-to-end category break would require stronger production
rank evidence or measurements and a substantially more complete
implementation-level cost account.

\paragraph{How to read the paper.}
A reader new to MQ can read \Cref{sec:mq-background} first, then focus on the
plain-language paragraphs before and after each theorem.  The exact attack is
in \Cref{sec:prelim,sec:symmetry,sec:structured-lift}.  Root existence and
retry probabilities are separated in \Cref{sec:slices}; the solver and gate
models are isolated in \Cref{sec:complexity}; and
\Cref{sec:evidence} states exactly which assumptions are tested rather than
proved.  All long proofs and exhaustive cost-search certificates are deferred
to the appendices.

The submitted $q=16$ matrices are not symmetrized and therefore do not satisfy
the hypothesis of our quotient theorem.  Returning unchanged to that design,
however, would restore the known characteristic-two attack surface discussed
in \Cref{sec:related}.

\section{An MQ Primer and Relation to Prior Work}\label{sec:mq-background}

This section supplies the minimum background needed for the attack.  Readers
already comfortable with underdetermined MQ, polar forms, Macaulay matrices,
and multigraded XL may skip to \Cref{sec:prelim}.

\subsection{MQ systems, roots, and affine slices}\label{sec:mq-primer}

Write $\F_q$ for the finite field with $q$ elements.  An MQ system over
$\F_q$ is a vector-valued polynomial map of total degree at most two,
\[
  g=(g_1,\ldots,g_K):\F_q^N\longrightarrow\F_q^K,
\]
together with a target $z\in\F_q^K$.  We continue to call $g$ quadratic when
some coordinates also contain linear or constant terms.  A \emph{root} is an
$x$ with $g(x)=z$.
The set
\[
  g^{-1}(z)=\{x:g(x)=z\}
\]
is the \emph{fiber} above $z$: it is the set of all candidate assignments
that produce the target.  After the affine reconstruction described later,
those assignments become full signature candidates.  All arithmetic is over
the stated finite field; every concrete \SNOVA{} calculation in this paper uses
$q=19$, so additions and multiplications are performed modulo $19$.

As a toy example, two equations such as
\[
 x_1x_2+x_3^2=z_1,
 \qquad
 x_1^2+x_2x_3+x_1=z_2
\]
form an MQ system in three variables.  Having more variables than equations
makes the system \emph{underdetermined}, but does not make it linear or
necessarily easy.  Modern MQ attacks often spend most of their time on very
large linear-algebra matrices built from the quadratic equations.

An \emph{affine space} is a translate $a+V=\{a+v:v\in V\}$ of a linear
subspace.  An \emph{affine slice} $a+L\subseteq a+V$ restricts the search to a
smaller direction space $L$.  If $h=\dim L$, choosing a basis of $L$ turns the
slice into an MQ system in $h$ coordinates.  Smaller $h$ makes the algebraic
solve cheaper, but it also lowers the chance that the slice contains a root;
\Cref{sec:slices} quantifies this tradeoff.

Three logically distinct events recur throughout the paper:
\begin{enumerate}[label=(\roman*),itemsep=2pt]
\item the chosen target and slice contain at least one $\F_q$-rational root;
\item the slice contains exactly one such root; and
\item the particular Macaulay matrix used by special XL has a one-dimensional
      right kernel from which row reduction reveals that root.
\end{enumerate}
The first two are properties of the quadratic map.  The third is a statement
about a specific solver representation.  A proof of root existence is not, by
itself, a proof that one solver invocation recovers the root.

\subsection{Rank, polar forms, and collisions}

For a scalar quadratic polynomial $f$, define its polar form by
\[
  \beta_f(x,y)=f(x+y)-f(x)-f(y)+f(0).
\]
In odd characteristic, $\beta_f$ is bilinear.  It is the finite-field analogue
of the part of a directional derivative that comes from the quadratic terms.
Its matrix rank measures how many independent directions the quadratic part
couples.

For a vector-valued map $g=(g_1,\ldots,g_K)$, every nonzero coefficient vector
$b\in\F_q^K$ selects a scalar combination $b\cdot g$.  The minimum polar rank
must be taken over \emph{all} such output directions.  Multiplying $b$ by a nonzero scalar does not change the rank, so these
vectors can be grouped into projective directions: one-dimensional lines
through the origin.

Why is rank useful?  Finite-field Fourier, or character-sum, identities imply
that a scalar quadratic with high polar rank is well balanced.  Applied to all nonzero combinations
$b\cdot g$, this gives rigorous bounds on every target fiber; when the ranks
are high enough relative to $K$, its size is close to the random-map prediction
$q^{N-K}$.  On a slice, a related
\emph{cross-polar rank} measures how strongly a displacement inside the slice
interacts with the full residual domain.  High cross-polar rank suppresses
pairs of distinct slice points that map to the same target.  Thus:
\begin{center}
\begin{tabularx}{0.91\textwidth}{@{}>{\bfseries}p{0.28\textwidth}X@{}}
Polar rank & controls whether full-space target fibers are nonempty and close
to their expected size.\\
Cross-polar rank & controls collisions inside an affine slice and therefore
unique-root retry probabilities.\\
Macaulay kernel & its dimension controls whether the chosen special-XL linear
algebra actually exposes the root.
\end{tabularx}
\end{center}
High rank is not invoked as a vague ``randomness'' heuristic in the structural
proofs: \Cref{sec:slices} gives exact first- and second-moment formulas.  The
production issue is whether the required rank lower bound holds in every
nonzero output direction.  Public sampling gives evidence, but only a proof or
exhaustive enumeration certifies the global minimum.

\subsection{XL and Macaulay matrices in one page}\label{sec:xl-primer}

XL (``eXtended Linearization'') turns a nonlinear system into a larger linear
system.  Starting from quadratic equations $f_i=0$, it multiplies each $f_i$
by selected monomials.  For example,
\[
  f=x_1x_2+x_3^2+x_1+1
  \quad\Longrightarrow\quad
  x_1f=x_1^2x_2+x_1x_3^2+x_1^2+x_1.
\]
Each resulting polynomial becomes a row.  Each monomial---for example
$x_1^2x_2$ or $x_1x_3^2$---becomes a column.  The coefficient matrix is a
\emph{Macaulay matrix}.  Gaussian elimination on this matrix can reveal enough
relations among the monomials to recover a root.  Treating monomials as columns
is only a linear-algebra device: the monomials are not independent variables in
the original problem, and the solver must eventually recover coordinates that
satisfy all of their multiplicative relations.

The cost comes from the size and rank of this matrix, not from evaluating the
original equations once.  Generic XL tracks only total degree.\@  \emph{Special
XL} exploits a variable-block structure.  A monomial then has a vector of
degrees, one entry per block, called its \emph{multidegree}.  Building only
selected multidegrees can make the Macaulay matrix much smaller.

Our $\ell=4$ reduction produces four blocks after a change to eigenblock
coordinates.  A quadratic using one block twice has degree pattern
$(2,0,0,0)$ up to permutation; a bilinear equation using two blocks has pattern
$(1,1,0,0)$.  The stable-lift construction in
\Cref{sec:structured-lift} exists precisely to preserve these patterns after
affine elimination.  A \emph{homogenizer} is one extra variable per block that
allows the constant and linear terms to be written with the same
multidegree; it is set to one when the root is read out.

If $x$ is a root, the vector of all selected monomial values
\[
  \nu(x)=(1,x_1,\ldots,x_N,x_1^2,x_1x_2,\ldots)
\]
lies in the right kernel of the dehomogenized Macaulay matrix.  The right
kernel is the space of column vectors $w$ with $\mathcal M w=0$.  If it is
one-dimensional, normalizing its constant coordinate to one reveals the
degree-one coordinates and hence the root.  A unique base-field root does
not automatically imply a one-dimensional Macaulay kernel: extension-field
roots or algebraic relations among rows can create extra kernel vectors.  This
is the separate production solver assumption in the structured attacks.

\subsection{Relation to prior attacks}\label{sec:related}

Beullens rewrote the \SNOVA{} verifier in ordered bilinear power-pair
coordinates and derived affine-column forgery attacks from rank defects in the
resulting outer matrix~\cite{Beullens2024SNOVA}.  Our quotient acts one layer
earlier: public-matrix symmetry identifies the power-pair polynomials
\emph{before} the outer map sees them.

Cabarcas et al.\ exploit stability under the public matrix algebra with a
multigraded XL algorithm~\cite{Cabarcas2025SNOVA}.  Furue et al.\ expose related
multihomogeneous systems through matrix transformations~\cite{FurueEtAl2025}.  Our attack starts from the complete affine forgery
system, uses a public column lift to retain every output equation, and then
restricts to an invariant kernel on which the affine offsets preserve the
block grading needed by special XL.\@

The wedge attack breaks the submitted characteristic-two parameters through
the block-ring structure~\cite{BrosEtAl2026SNOVAWedge}; later work studies the
wedge map further~\cite{TsengEtAl2026Wedge}.  Other analyses reinterpret
scalar-expanded \SNOVA{} as Unbalanced Oil and Vinegar (UOV), or lift it to
extension-field UOV, for key recovery~%
\cite{IkematsuAkiyama2024,LiDing2024,NakamuraEtAl2025}.  Symmetric and
exterior algebras also appear in recent UOV key-recovery attacks~%
\cite{JinEtAl2026,Ran2026}.  Our attack is a direct forgery attack: it does not
recover the hidden oil subspace that constitutes the UOV-style secret key.

For the generic $\ell=2$ cores, we reproduce Hashimoto's optimization,
including its published boundary value $a=1$~%
\cite{Hashimoto2023,Beullens2024SNOVA}.  When this leaves the elementary system
$\operatorname{MQ}(q,1,1)$, the artifact solves it by exhaustive evaluation and
charges that work conservatively.  May, Ostuzzi, and Ressler's \emph{Just
Guess} algorithm and the generalized variable partitioning of Asanuma et al.\ are more recent approaches to underdetermined MQ~%
\cite{MayOstuzziRessler2026,AsanumaEtAl2026}.  Evaluating the released
classical Just-Guess formulas on our residual dimensions gives larger estimates
than the selected Beullens--Hashimoto modes.  We have not implemented the newer
generalized partition search.  The ``generic MQ'' numbers below are therefore
pinned, reproducible baselines rather than a claim that every possible generic
solver has been exhausted.

Solver engineering is also moving quickly.  Sakata and Takagi use
Hilbert-series information to reduce the matrices built by an $F_4$
Gr\"obner-basis solver and report a new MQ-challenge
record~\cite{SakataTakagi2026}.  That work is an
implementation and challenge result rather than a plug-in cost formula for the
underdetermined residual systems here, so we do not substitute it into
\Cref{tab:main}.  It reinforces the need to distinguish exact algebraic
reductions, solver-operation models, and complete implementations.


\subsection{Notation at a glance}\label{sec:notation}

\begin{table}[H]
\centering
\caption{Recurring notation.  Generic primer notation such as $N$ is local;
$N_0$ is the specific common-column dimension in the \SNOVA{} reduction.}
\label{tab:notation}
\small
\renewcommand{\arraystretch}{1.12}
\begin{tabularx}{0.96\textwidth}{@{}>{\raggedright\arraybackslash}p{0.23\textwidth}X@{}}
\toprule
Symbol & Meaning\\
\midrule
$(v,o,\ell,r)$ & Public parameter tuple: $v$ vinegar blocks, $o$ oil blocks,
block height $\ell$, and $r$ columns per signature block.\\
$n=v+o$ & Number of matrix-valued signature blocks.\\
$M=or\ell$ & Number of scalar public output equations.\\
$N$ & Generic domain dimension for an MQ map $g:\F_q^N\to\F_q^K$; its
meaning is local to the section in which the map is defined.\\
$N_0=n\ell$ & Length of one stacked signature column after the common-column
relation.\\
$N_{\rm raw}=m_1\ell^2r^2$ & Number of labeled verifier feature slots before imposing
the relation.\\
$A=\F_q[S]$ & Public commutative matrix algebra generated by $S$; in the
Version~2.3 rows, $\dim_{\F_q}A=\ell$.\\
$m_1$ & Number of public base-form groups in the pinned implementation.\\
$K=m_1\binom{\ell+1}{2}$ & Maximum number of distinct quadratic power-pair
features after the symmetry quotient.\\
$Q_R$ & Matrix mapping the $K$ quotient features to public output directions.\\
$C_{R,\mathbf v}$ & Linear constraint matrix obtained by projecting away the quadratic
image.\\
$H_{R,\mathbf v}$ & Stable $A$-linear kernel used to choose structured solver slices.\\
$L$, $h$ & Slice direction space and its base-field dimension.\\
$\tau=K-h$ & Slice deficit.  It enters the exact full-rank retry formula;
for the selected $\tau=10,14$ slices, the retry count is essentially $q^\tau$.\\
$r_*$ & Minimum full-space polar rank over all nonzero scalar output
combinations.\\
$e_L(b)$ & Cross-polar rank of output direction $b$ relative to slice $L$.\\
$\mu$, $\bar\eta$ & Expected roots on a random target--slice pair and an upper
bound on its collision mass.\\
\bottomrule
\end{tabularx}
\end{table}

\section{The Affine-Column Framework}\label{sec:prelim}

This section rewrites the public verifier after an attacker-imposed affine
relation among signature columns.  The result has three visibly separate
parts: quotient quadratic features, linear terms created by the offsets, and a
constant term.  \Cref{prop:affine-reduction} then eliminates the affine part
exactly.

Three vector spaces appear and should not be confused.  The raw public feature
vector has length $N_{\rm raw}$; the verifier outputs $M$ scalar equations; and the
common column variable $u$ has length $N_0$.  The matrices below simply move
between these three spaces.

\subsection{Parameters and public forms}\label{sec:parameters}

The public parameter tuple is $(v,o,\ell,r)$.  There are $v$ vinegar blocks and
$o$ oil blocks, for a total of
\[
  n=v+o
\]
matrix-valued signature blocks.  The names ``vinegar'' and ``oil'' come from
SNOVA's UOV ancestry; in this public forgery attack they serve only to specify
dimensions, and the attacker never learns or reconstructs the secret partition.
Each block has $\ell$ rows and $r$ columns.
Stacking the $j$th column of all $n$ blocks gives a vector of length
$N_0=n\ell$; the verifier has $M=or\ell$ scalar outputs.  The pinned
Version~2.3 implementation defines
\begin{equation}\label{eq:parameters-basic}
 m_1=\left\lceil\frac{or}{\ell}\right\rceil,
 \qquad M=or\ell,
 \qquad N_0=n\ell.
\end{equation}
Here $m_1$ is the number of public base-form groups used to assemble the
outputs.  The ceiling is used by the reference implementation; a public
analysis script uses a floor and therefore disagrees when $\ell\nmid or$.  All
dimensions in this paper follow the pinned implementation~\cite{SNOVA23}.

Let $S\in\Mat_\ell(\F_q)$ be the public symmetric matrix with irreducible
characteristic polynomial.  Its powers generate the commutative matrix
algebra
\[
 A=\F_q[S],\qquad \Sigma_\xi=I_n\otimes\xi\quad(\xi\in A).
\]
For the Version~2.3 parameters, $A\cong\F_{q^\ell}$ and
$\dim_{\F_q}A=\ell$.  Multiplication by $\Sigma_\xi$ applies the same
$\ell\times\ell$ matrix $\xi$ independently to every signature block.

The scalar-expanded public base matrices are
$P_1,\ldots,P_{m_1}\in\Mat_{N_0}(\F_q)$.  In the public $q=19$ construction,
each $P_i$ is symmetric.  For $\xi,\eta\in A$, define the bilinear form
\begin{equation}\label{eq:bilinear-form}
 B_i^{\xi,\eta}(x,y)=x^\transpose\Sigma_\xi P_i\Sigma_\eta y.
\end{equation}
Setting $x=y$ makes this a quadratic feature.  The power basis
$1,S,\ldots,S^{\ell-1}$ gives the ordered power-pair coordinates used in the
prior affine-column attack.

\subsection{Imposing one common column variable}\label{sec:affine-columns}

Write the stacked signature columns as
\[
  U=[u_0|\cdots|u_{r-1}],\qquad u_j\in\F_q^{N_0}.
\]
The attacker chooses multipliers $R_j\in A$ and offsets
$v_j\in\F_q^{N_0}$, with $R_0=1$.  Write
$R=(R_0,\ldots,R_{r-1})$ and
$\mathbf v=(v_0,\ldots,v_{r-1})$, and impose
\begin{equation}\label{eq:relation}
 u_j=\Sigma_{R_j}u+v_j,
 \qquad 0\le j<r.
\end{equation}
Thus all $r$ columns are affine functions of one vector $u$.  This is not an
assumption about honest signatures; it defines a publicly checkable subset in
which the attacker searches for a forgery.

Before the relation is imposed, the feature vector has
\[
  N_{\rm raw}=m_1\ell^2r^2
\]
coordinates: $m_1$ base forms, $\ell^2$ ordered power pairs, and $r^2$ ordered
column pairs.  The official verifier flattens each $r\times\ell$ block in
$(i,j)$ order, whereas our feature formulas use the block-transposed $(j,i)$
order.  Let $t_{\rm off}\in\F_q^M$ and
$\mathcal E_{\rm off}\in\F_q^{M\times N_{\rm raw}}$ denote the target and public outer
map in the official order.  If $\Pi_{\rm out}$ is the corresponding public
permutation, define
\[
  t=\Pi_{\rm out}t_{\rm off},\qquad
  \mathcal E=\Pi_{\rm out}\mathcal E_{\rm off}.
\]
This is only a coordinate reordering; it changes neither ranks nor solutions.

After substituting \eqref{eq:relation}, products of two variable parts give
quadratic terms in $u$, products of one variable part and one offset give
linear terms, and products of two offsets give constants.  For a fixed
relation $R=(R_j)_{j=0}^{r-1}$, let
$\mathcal R_R\in\F_q^{N_{\rm raw}\times m_1\ell^2}$ expand the remaining ordered
quadratic features into the full feature vector, and let
$\mathcal F_{R,\mathbf v}\in\F_q^{N_{\rm raw}\times N_0}$ collect the offset-linear
coefficients.  The complete substituted verifier is
\begin{equation}\label{eq:ordered-system}
 \mathcal E\bigl(\mathcal R_Rz_{\rm ord}(u)+\mathcal F_{R,\mathbf v}u\bigr)
 +c_{R,\mathbf v}=t,
\end{equation}
where $z_{\rm ord}(u)$ contains the $m_1\ell^2$ values
$B_i^{S^a,S^b}(u,u)$ and $c_{R,\mathbf v}$ is the public constant term.  Put
\begin{equation}\label{eq:ER}
 E_R=\mathcal E\mathcal R_R\in\F_q^{M\times m_1\ell^2},
 \qquad \Lambda_{R,\mathbf v}=\mathcal E\mathcal F_{R,\mathbf v}\in\F_q^{M\times N_0}.
\end{equation}

\subsection{Separating quadratic and affine output directions}

By \Cref{thm:symmetric-square}, the ordered feature vector duplicates every
off-diagonal power pair.  Let $z_{\rm sym}(u)$ contain one coordinate for each
unordered pair $0\le a\le b<\ell$, and let $D$ copy each off-diagonal
coordinate into its two ordered positions.  Then
$z_{\rm ord}(u)=Dz_{\rm sym}(u)$.  Define
\begin{equation}\label{eq:QR}
 Q_R=E_RD\in\F_q^{M\times K},
 \qquad K=m_1\binom{\ell+1}{2}.
\end{equation}
The columns of $Q_R$ span all output directions that the residual quadratic
features can reach.

Now choose $W_R$ whose rows form a basis of the left kernel of $Q_R$ and put
\begin{equation}\label{eq:C}
 C_{R,\mathbf v}=W_R\Lambda_{R,\mathbf v}.
\end{equation}
Multiplying by $W_R$ kills every quadratic feature.  What remains is therefore
an ordinary affine system in $u$.  This elementary projection is the key to
retaining the \emph{complete} verifier rather than dropping output equations
that do not lie in the quadratic image.

\begin{center}
\begin{tabularx}{0.92\textwidth}{@{}>{\bfseries}p{0.18\textwidth}X@{}}
$Q_R$ & sends the $K$ distinct quadratic features to the $M$ public output
coordinates; its column space is the quadratic image.\\
$W_R$ & keeps precisely the output directions orthogonal to that image, so
$W_R Q_R=0$.\\
$C_{R,\mathbf v}=W_R\Lambda_{R,\mathbf v}$ & records the ordinary linear equations that remain in
those output directions after the offsets are introduced.
\end{tabularx}
\end{center}

\begin{center}
\fcolorbox{deepblue}{lightblue}{\parbox{0.92\textwidth}{\small
\textbf{A two-output example.}  Suppose that, after the column substitution,
two verifier outputs over $\F_{19}$ are
\[
  y_1=q(u)+\ell_1(u),\qquad y_2=2q(u)+\ell_2(u),
\]
where $q$ is quadratic and $\ell_1,\ell_2$ are affine.  The combination
$2y_1-y_2=2\ell_1(u)-\ell_2(u)$ cancels the quadratic term.  Solving this
affine equation and substituting it back leaves one quadratic equation; neither
original output was dropped.  The rows of $W_R$ perform the same cancellation
simultaneously for the complete verifier.}}
\end{center}

\begin{proposition}[Exact affine reduction]\label{prop:affine-reduction}
Let $\rho=\rank Q_R$, $\lambda=\rank C_{R,\mathbf v}$, and
$\delta=M-\rho-\lambda$.  Then
\[
 \rank[Q_R\ \Lambda_{R,\mathbf v}]=\rho+\lambda,
\]
where $[Q_R\ \Lambda_{R,\mathbf v}]$ denotes horizontal matrix concatenation.  After row
reduction, exactly $\delta$ independent affine consistency conditions remain
on the target.  For every target satisfying those conditions, Gaussian
elimination leaves $\rho$ equations of degree at most two in
$N_0-\lambda$ variables, with exactly the same solution set as
\eqref{eq:ordered-system}.

In particular, if $\rho=K$ and $\lambda=M-K$, every target passes the
\emph{affine consistency test}, and the residual system has $K$ equations of
degree at most two on an affine translate of $\ker C_{R,\mathbf v}$.  This conclusion
does not assert that the residual system has a root.
\end{proposition}

\begin{proof}
The row space of $W_R$ is $\ker Q_R^\transpose$.  Since $W_R$ has full row
rank,
\[
 \ker W_R=\bigl(\ker Q_R^\transpose\bigr)^\perp=\operatorname{im}Q_R.
\]
Thus $W_R$ realizes the quotient map
$\F_q^M\to\F_q^M/\operatorname{im}Q_R$.  The image of the columns of
$\Lambda_{R,\mathbf v}$ in this quotient has dimension
$\rank(W_R \Lambda_{R,\mathbf v})=\lambda$, which proves
$\rank[Q_R\ \Lambda_{R,\mathbf v}]=\rho+\lambda$.

Write $b=t-c_{R,\mathbf v}$.  Projecting the substituted verifier by $W_R$ gives
\[
 C_{R,\mathbf v}u=W_R b.
\]
The codomain of $C_{R,\mathbf v}$ has dimension $M-\rho$, while its image has
dimension $\lambda$.  Solvability therefore imposes exactly
$(M-\rho)-\lambda=\delta$ independent affine conditions on the target.
For a target satisfying the affine consistency conditions, parameterize the
solution space of these affine equations by $N_0-\lambda$ free variables and
substitute into a complementary set of $\rho$ rows.  This leaves $\rho$
equations of degree at most two and uses only invertible row operations and
exact affine elimination, so the solution set is unchanged.  If $\rho=K$ and
$\lambda=M-K$, then $\delta=0$ and $C_{R,\mathbf v}$ is surjective onto its
$M-K$-dimensional codomain; hence every target passes the affine consistency
test.
\end{proof}

\paragraph{Operational meaning.}
The attacker computes $Q_R$, $W_R$, and $C_{R,\mathbf v}$ from the public key before
running an MQ solver.  If the ranks are unfavorable, the relation or offsets
are changed.  When $\rank Q_R=K$ and $\rank C_{R,\mathbf v}=M-K$, every hash target first
determines a nonempty affine space $a(t)+\ker C_{R,\mathbf v}$ on which
$K$ equations of
degree at most two remain.  This is an exact reformulation, not a root-existence
theorem: the residual equations may still have no solution.  Root availability
is analyzed separately in \Cref{sec:slices}.  No verifier equation has been
discarded.  Conversely, every root of the residual system maps through the
chosen column relation to a solution of the complete substituted verifier.

\section{The Symmetry Quotient}\label{sec:symmetry}

The central reduction is the formal version of
\eqref{eq:intro-swap}.  Readers may interpret $\Sym^2(A)$ simply as the vector
space spanned by \emph{unordered} pairs from $A$; the tensor notation records
why the identification is basis-independent.

For a fixed base form $P_i$, define
\[
  q_{i,a,b}(u)=u^\transpose\Sigma_{S^a}P_i\Sigma_{S^b}u.
\]
Transposition gives $q_{i,a,b}=q_{i,b,a}$.  In coordinates, the feature vector
therefore contains two copies of every off-diagonal polynomial.  The theorem
below shows that this is not an accident of the power basis: every alternating
label $\xi\otimes\eta-\eta\otimes\xi$ vanishes.  Write
$\F_q[u]^{\mathrm{hom}}_2$ for the vector space of homogeneous quadratic
polynomials in the coordinates of $u$.

\begin{theorem}[Symmetric-square factorization]\label{thm:symmetric-square}
Let $A=\F_q[S]$ have dimension $d$, and assume $S^\transpose=S$ and
$P_i^\transpose=P_i$.  For each $i$, the map
\[
 A\otimes_{\F_q}A\longrightarrow \F_q[u]^{\mathrm{hom}}_2,
 \qquad
 \xi\otimes\eta\longmapsto
 u^\transpose\Sigma_\xi P_i\Sigma_\eta u
\]
is invariant under $\xi\otimes\eta\mapsto\eta\otimes\xi$.  It therefore
factors through
\[
 \Sym^2_{\F_q}(A)
 =\bigl(A\otimes_{\F_q}A\bigr)
   \big/\Span\{\xi\otimes\eta-\eta\otimes\xi\}.
\]
Consequently, the direct sum over the $m_1$ base forms has rank at most
$m_1\binom{d+1}{2}$.  After any public linear map to $M$ outputs---in
particular, for every relation-dependent matrix $Q_R$---the rank is at most
\begin{equation}\label{eq:symmetry-rank}
 \min\left\{M,m_1\binom{d+1}{2}\right\}.
\end{equation}
If the minimal polynomial of $S$ has degree $\ell$, then $d=\ell$ and, in the
power basis,
\begin{equation}\label{eq:duplication}
 z_{\rm ord}(u)=Dz_{\rm sym}(u).
\end{equation}
\end{theorem}

\begin{proof}
Every element of $A$ is symmetric.  Since the displayed expression is a
scalar,
\[
 u^\transpose\Sigma_\xi P_i\Sigma_\eta u
 =\bigl(u^\transpose\Sigma_\xi P_i\Sigma_\eta u\bigr)^\transpose
 =u^\transpose\Sigma_\eta P_i\Sigma_\xi u.
\]
Thus every alternating label $\xi\otimes\eta-\eta\otimes\xi$ lies in the
kernel.  Taking the direct sum over $i$ gives the first rank bound, and a
linear map to $M$ public outputs gives the additional codomain cap.  When the
powers of $S$ form a basis of $A$, the factorization is exactly the coordinate
identity \eqref{eq:duplication}.
\end{proof}

\paragraph{Why the outer map cannot repair the loss.}
The public outer map receives only the values of these quadratic features.  If
two labeled inputs already define the same polynomial, no later linear
combination can separate them.  The outer map may have full rank on the
$K$-dimensional quotient, but it cannot enlarge that quotient.

For $\ell=4$, each base form loses six of its sixteen ordered coordinates:
$\ell^2=16$ becomes $\binom{\ell+1}{2}=10$.  For $\ell=2$, four ordered
coordinates become three.  These reductions occur for every relation and every
symmetric public key.

The theorem is an upper bound: additional public dependencies could make the
quadratic image even smaller.  The concrete attacks separately check
$\rank Q_R=K$ before solving.  Thus the value $K$ used in the cost tables is
not inferred from symmetry alone; it is the exact rank observed in the public
preflight for the chosen relation.

The factorization itself uses only transposition symmetry and is valid in every
characteristic.  Odd characteristic is needed only for the familiar direct-sum
interpretation
$A\otimes A=\Sym^2(A)\oplus\bigwedge^2(A)$.  The submitted $q=16$ system is
outside this theorem because its public base matrices are not symmetrized.

\section{Structured Stable Lifts}\label{sec:structured-lift}

\Cref{prop:affine-reduction} permits arbitrary offsets $v_j$.  Generic offsets
are useful for making the affine equations full rank, but they create linear
terms that mix all coordinates.  That mixing destroys the block grading used
by the selected published $\ell=4$ special-XL model.  We therefore need offsets
that do two jobs at once:
\begin{enumerate}[label=(\roman*),itemsep=2pt]
\item span the missing affine output directions, so that the full verifier is
      retained; and
\item confine all offset-linear terms to a small public row space, so that a
      large invariant kernel remains available for special XL.\@
\end{enumerate}

The following construction achieves both.  Choose one vector
$w\in\F_q^{N_0}$ and multipliers $\Gamma_0,\ldots,\Gamma_{r-1}\in A$, and set
\begin{equation}\label{eq:structured-offsets}
 v_j=\Sigma_{\Gamma_j}w.
\end{equation}
All offsets are therefore generated from the same $w$ by the public algebra
$A$.  Define the row space
\begin{equation}\label{eq:Uw}
 \mathcal U_w=
 \Span_{\F_q}\left\{
 w^\transpose\Sigma_\xi P_i\Sigma_\eta:
 1\le i\le m_1,\ \xi,\eta\in A
 \right\}
 \subseteq(\F_q^{N_0})^*.
\end{equation}
This space contains every linear form that can arise from crossing a variable
term with one of the structured offsets.  Its a priori bound is $m_1d^2$, not
$m_1\binom{d+1}{2}$: one argument is the fixed offset vector $w$ rather than
the variable $u$ on both sides, so reversing the two algebra labels need not
leave this \emph{linear} form unchanged.

A subspace $H\subseteq\F_q^{N_0}$ is called $A$-linear if
$\Sigma_\xi H\subseteq H$ for every $\xi\in A$.  Since
$A\cong\F_{q^\ell}$, this is the analogue of extension-field linearity.  To
force such a kernel, take the complete $A$-orbit of the affine constraints.
For a basis $1,S,\ldots,S^{d-1}$ of $A$, form
\begin{equation}\label{eq:orbit-matrix}
 \mathcal O_{R,\mathbf v}=
 \begin{bmatrix}
 C_{R,\mathbf v}\\ C_{R,\mathbf v}\Sigma_S\\ \vdots\\ C_{R,\mathbf v}\Sigma_{S^{d-1}}
 \end{bmatrix}.
\end{equation}

\begin{theorem}[Stable-lift kernel]\label{thm:stable-kernel}
Retain the assumptions $S^\transpose=S$ and $P_i^\transpose=P_i$ from
\Cref{thm:symmetric-square}.  For offsets of the form
\eqref{eq:structured-offsets}:
\begin{enumerate}[label=(\roman*)]
\item $\mathcal U_w$ is closed under right multiplication by $A$ and
$\dim_{\F_q}\mathcal U_w\le m_1d^2$;
\item every row of $\mathcal F_{R,\mathbf v}$ and $\mathcal O_{R,\mathbf v}$ lies in
$\mathcal U_w$;
\item if $\rank\mathcal O_{R,\mathbf v}=m_1d^2$, then
\[
 H_{R,\mathbf v}=\ker\mathcal O_{R,\mathbf v}
\]
is $A$-linear, has base-field dimension $N_0-m_1d^2$, and satisfies
$\mathcal F_{R,\mathbf v}H_{R,\mathbf v}=0$.
\end{enumerate}
For the irreducible degree-$\ell$ construction, $H_{R,\mathbf v}$ has $A$-dimension
$n-m_1\ell$.
\end{theorem}

\begin{proof}
See \Cref{app:proof-stable-kernel}.
\end{proof}

\paragraph{Interpretation.}
When the public orbit matrix has the stated rank, its kernel is a large search
direction space closed under $A$.  Moving inside this kernel changes the
quadratic features but changes none of the offset-linear features.  The
offsets can therefore solve the missing affine equations, while the nonlinear
solver sees a clean $A$-linear slice.

\begin{center}
\fcolorbox{deepblue}{lightblue}{\parbox{0.94\textwidth}{\small
\textbf{The stable-lift logic in four arrows.}
Structured offsets force every offset-linear row into $\mathcal U_w$;
full orbit rank makes $\operatorname{row}(\mathcal O_{R,\mathbf v})=\mathcal U_w$;
therefore $H_{R,\mathbf v}=\ker\mathcal O_{R,\mathbf v}$ is annihilated by
$\mathcal F_{R,\mathbf v}$; and the complete $A$-orbit in $\mathcal O_{R,\mathbf v}$ makes
$H_{R,\mathbf v}$ $A$-linear.  The offsets can thus cover the missing affine output
directions without destroying the block structure used by special XL.\@}}
\end{center}

\subsection{From an \texorpdfstring{$A$}{A}-linear kernel to four variable blocks}

Special XL needs variables divided into blocks with controlled
multidegrees.  The algebra $A$ supplies those blocks after extending scalars to
a field $\Omega$ in which the minimal polynomial of $S$ splits.  Over
$\Omega$, multiplication by $S$ diagonalizes into $\ell$ eigenspaces.  The
primitive idempotent $e_a$ is simply the projection onto the $a$th eigenspace.
This scalar extension is a solver coordinate change, not a relaxation of the
signature problem: the recovered point must still descend to $\F_q$ and pass
the original verifier.

Let $\mathbf e_a\in\mathbb Z^\ell$ denote the $a$th standard unit vector.  A
quadratic using only block $a$ has multidegree $2\mathbf e_a$; a bilinear term
using blocks $a$ and $b$ has multidegree
$\mathbf e_a+\mathbf e_b$.

\begin{lemma}[Eigenblock coordinates of the quotient]\label{lem:spectral-quotient}
Let $\Omega/\F_q$ be a splitting field of the minimal polynomial of $S$, and
let $e_1,\ldots,e_\ell$ be the primitive idempotents of
$A\otimes_{\F_q}\Omega$.  If $L$ is an $A$-linear space of $A$-dimension
$s$, then
\[
 L\otimes_{\F_q}\Omega=L_1\oplus\cdots\oplus L_\ell,
 \qquad L_a=\Sigma_{e_a}(L\otimes_{\F_q}\Omega),
 \qquad \dim_\Omega L_a=s.
\]
For each $i$, the eigenblock features
\[
 x^\transpose\Sigma_{e_a}P_i\Sigma_{e_b}x,
 \qquad 1\le a\le b\le\ell,
\]
are an invertible linear recombination of the symmetric power-pair features.
They have multidegree $2\mathbf e_a$ when $a=b$ and
$\mathbf e_a+\mathbf e_b$ when $a<b$.
\end{lemma}

\begin{proof}
Finite fields are perfect, so the irreducible minimal polynomial of $S$ is
separable.  Over $\Omega$, the power basis and the primitive-idempotent basis
of $A\otimes\Omega\cong\Omega^\ell$ are related by an invertible Vandermonde
matrix.  Passing to the symmetric square gives an invertible change between the
unordered power-pair features and the unordered eigenblock features.  The
idempotent $e_a$ projects $L\otimes\Omega$ onto $L_a$.  Hence a diagonal
feature uses only variables in $L_a$, while an off-diagonal feature is bilinear
in the variables from $L_a$ and $L_b$, giving the stated multidegrees.
\end{proof}

\begin{corollary}[Stable affine slices]\label{cor:stable-slices}
Assume the public checks
\begin{equation}\label{eq:stable-preflight}
 \rank Q_R=K,\qquad
 \rank C_{R,\mathbf v}=M-K,\qquad
 \rank\mathcal O_{R,\mathbf v}=m_1\ell^2
\end{equation}
pass, and let $L\subseteq H_{R,\mathbf v}$ be an $A$-linear subspace of
$A$-dimension $s$ (equivalently, base-field dimension $h=\ell s$).  For every
target, the affine equations in \Cref{prop:affine-reduction} have a
solution space $a(t)+\ker C_{R,\mathbf v}$ for a publicly computed base point
$a(t)$.  For every $y\in\ker C_{R,\mathbf v}$, the affine
slice $a(t)+y+L$ is contained in that space, and the full offset-linear feature
$\mathcal F_{R,\mathbf v}u$ is constant on the slice.

After extension to a splitting field, use the eigenblock coordinates of
\Cref{lem:spectral-quotient}.  Translation by $a(t)+y$ creates only linear and
constant terms in the same one or two blocks as the corresponding quadratic
term.  Introducing one homogenizing variable per block---an auxiliary variable that
is set to one after solving---therefore turns each complete affine equation
into a homogeneous equation of multidegree
$2\mathbf e_a$ or $\mathbf e_a+\mathbf e_b$.  If the restricted quadratic
feature map has total rank $K$, then the disjoint monomial supports of these
unordered multidegrees force rank $m_1$ in every component: there are
$\binom{\ell+1}{2}$ components, each has rank at most $m_1$, and their ranks
sum to $K=m_1\binom{\ell+1}{2}$.
\end{corollary}

\begin{proof}
See \Cref{app:proof-stable-slices}.
\end{proof}

\paragraph{What the preflight certifies.}
The first two ranks in \eqref{eq:stable-preflight} say that the quotient has
its full expected rank and that the offsets supply every remaining affine
output direction.  The orbit rank says that the kernel has the predicted
$A$-linear dimension and kills all offset-linear motion.  The restricted
quadratic rank then certifies that no multidegree component was lost on the
chosen slice.  All of these tests are public and occur before the expensive
Macaulay computation.  A failing relation, offset family, or slice is simply
replaced.

\subsection{Passing the verifier's block-symmetry check}\label{sec:rejection}

The public verifier does more than evaluate the quadratic map.  It also rejects
signatures containing too many symmetric blocks.  A root of the reduced MQ
system is useful only if it passes this final check.  We therefore build
rejection compatibility into the reduction rather than treating it as an
unmodeled success probability.

Version~2.3 rejects a square odd-characteristic signature if one of its blocks
is symmetric when $\ell>2$, and rejects an $\ell=2$ signature if more than
$\lfloor n/4\rfloor$ blocks are symmetric~\cite{SNOVA23}.  Rectangular blocks
are not subject to this rule.

\paragraph{The structured $\ell=4$ rows.}
For each square block $b$, substituting \eqref{eq:relation} into the equalities
$U_b[a,c]=U_b[c,a]$ gives an affine system
\begin{equation}\label{eq:rejection-affine}
 G_bu=-h_b.
\end{equation}
If $\rank[G_b\mid-h_b]>\rank G_b$, this system is inconsistent, so no point in
the entire residual search space can make block $b$ symmetric.  Requiring this
inconsistency for every square block is an exact public certificate that every
root produced by the structured lift passes rejection.

\paragraph{A deterministic $\ell=2$ construction.}
The first two parameter levels use offsets that force a fixed nonzero skew in
every block.

\begin{proposition}[Fixed-skew $\ell=2$ lift]\label{prop:fixed-skew}
For the Version~2.3 $\ell=2$ matrix
\[
 S=\begin{pmatrix}1&2\\2&15\end{pmatrix},
 \qquad R_0=I,
 \qquad R_1=9I+10S=\begin{pmatrix}0&1\\1&7\end{pmatrix},
\]
choose offsets satisfying
\begin{equation}\label{eq:fixed-skew-offset}
 (v_1)_{b,0}-(v_0)_{b,1}=1
\end{equation}
for every block $b$.  Then every block of every signature satisfying the
column relation has
\begin{equation}\label{eq:fixed-skew-identity}
 U_b[0,1]-U_b[1,0]=1.
\end{equation}
Consequently every solution passes the $\ell=2$ symmetric-block rejection
check.  If the public preflight also gives
$\rank Q_R=K$ and $\rank C_{R,\mathbf v}=M-K$, every target passes the
affine consistency test and leaves $K$ equations of degree at most two in
$N_0-M+K$ variables.  Every root of that residual system is a valid signature.
\end{proposition}

\begin{proof}
The first row of $R_1$ is the second row of $R_0=I$.  Hence the variable part
of $U_b[0,1]-U_b[1,0]$ vanishes.  The remaining offset difference is one by
\eqref{eq:fixed-skew-offset}, proving \eqref{eq:fixed-skew-identity}.  No block
is symmetric.  The final statement is the full-lift case of
\Cref{prop:affine-reduction}.
\end{proof}

Thus the fixed-skew construction has no rejection probability: every MQ root
is acceptable.  The Level-V headline attack uses a faster alternative in which
targets are filtered before solving.

\paragraph{The zero-offset $\ell=2$ filter.}
Set $R_0=I$, $R_1=0$, and use zero offsets.  Each block then has the form
\[
 \begin{pmatrix}x&0\\ y&0\end{pmatrix},
\]
so it is symmetric exactly when $y=0$.  These rejection events are nonzero
hyperplanes on disjoint block coordinates.  This relation is chosen
carefully: a relation that forced the two off-diagonal entries to agree would
make every block symmetric and would be useless.

Let $N$ be the residual-domain dimension ($N=N_0$ for this zero-offset lift)
and let $g:\F_q^N\to\F_q^K$ be the residual map of degree at most two.
For each nonzero
$b\in\F_q^K$, take the polar rank of the scalar quadratic $b\cdot g$, and let
$r_*$ be the minimum of these ranks.  Thus $r_*$ measures the worst output
direction; a large $r_*$ forces every target fiber to have nearly its expected
size.

Put
\[
 t_{\rm rej}=\lfloor n/4\rfloor+1,
 \qquad B=\binom n{t_{\rm rej}},
 \qquad \alpha=1-q^{-K},
\]
and let $\mathcal X_{\rm acc}$ be the assignments accepted by the block-symmetry filter.
At least $t_{\rm rej}$ symmetric blocks are needed for rejection.  Because
the $n$ hyperplanes use disjoint coordinates, the exact density of accepted
assignments is
\begin{equation}\label{eq:l2-accepted-density}
 a_n\coloneqq\frac{|\mathcal X_{\rm acc}|}{q^N}
 =q^{-n}\sum_{k=0}^{\lfloor n/4\rfloor}\binom nk(q-1)^{n-k}.
\end{equation}

\begin{proposition}[Accepted targets for the $\ell=2$ filter]\label{prop:l2-rejection}
Every target has an accepted root if
\begin{equation}\label{eq:l2-rejection}
 1-Bq^{-t_{\rm rej}}>\alpha(1+B)q^{K-r_*/2}.
\end{equation}
For a uniform target $z\in\F_q^K$,
\begin{equation}\label{eq:l2-accepted-target}
 \Pr\bigl[g^{-1}(z)\cap\mathcal X_{\rm acc}\ne\varnothing\bigr]
 \ge \frac{a_n}{1+\alpha q^{K-r_*/2}}.
\end{equation}
Moreover, if $1-\alpha q^{K-r_*/2}>0$, every target $z$ satisfies
\begin{equation}\label{eq:l2-accepted-fraction}
 \frac{|g^{-1}(z)\cap\mathcal X_{\rm acc}|}{|g^{-1}(z)|}
 \ge
 1-\frac{Bq^{-t_{\rm rej}}+\alpha Bq^{K-r_*/2}}
          {1-\alpha q^{K-r_*/2}},
\end{equation}
whenever the right-hand side is positive.
\end{proposition}

\begin{proof}
See \Cref{app:proof-l2-filter}.
\end{proof}

The three conclusions answer different questions.  Readers interested only
in the concrete attack may retain the following takeaway: for the Level-V
system, the global condition $r_*\ge240$ implies the first inequality and
therefore gives an accepted root for every target.  More generally,
\eqref{eq:l2-rejection} is a sufficient worst-case condition: every target has
an accepted root.  Inequality \eqref{eq:l2-accepted-target} lower-bounds the
fraction of uniformly sampled targets that have one.  Inequality
\eqref{eq:l2-accepted-fraction} says that, target by target, almost every root
is accepted when the rank is high enough.  The Level-V headline uses the first
condition; the second gives an explicit retry factor if only a weaker global
rank bound is available.

\section{Root Existence, Uniqueness, and Retry Bounds}\label{sec:slices}

After the exact affine reduction, the structured attack repeatedly chooses a
hash target and an affine coset of the selected slice direction space.  We call
one such pair a \emph{target--slice trial}.  The attack asks whether the
resulting MQ system has a recoverable root.  Three events must not be conflated:
\begin{enumerate}[label=(\roman*),itemsep=2pt]
\item the slice contains at least one base-field root;
\item it contains exactly one base-field root; and
\item the chosen Macaulay matrix has a one-dimensional kernel from which that
      root can be read.
\end{enumerate}
This section proves statements about the first two events.  The third is the
separate special-XL solving assumption in \Cref{sec:special-xl}.

Let $V$ be the full residual direction space, let $L\subseteq V$ be the chosen
slice direction space, and write $h=\dim L$.  For a map of degree at most two
$g:V\to\F_q^K$, a random-map heuristic predicts
\[
  \mu=q^{h-K}
\]
roots in a random affine coset of $L$ above a random target.  When $h=K-\tau$ with $\tau\ge0$, this mean is $q^{-\tau}$, suggesting
about $q^\tau$ trials.  The obstruction detected by a second-moment analysis is excess
collisions: distinct points in the same slice mapping to the same target.
Polar derivatives count those collisions exactly, so the argument below is a
theorem rather than a random-system heuristic.  The cleanest case is full
cross-polar rank: then distinct points behave almost as independently as the
mean-root calculation suggests.  For the positive slice deficits used here, a
unique-root trial costs essentially $q^{K-h}=q^\tau$ attempts.  The theorem first states the exact formulas and
then quantifies what changes when some output directions lose rank.

Write the vector-valued polar map as
\[
 \beta(x,y)=g(x+y)-g(x)-g(y)+g(0).
\]
For $b\in\F_q^K\setminus\{0\}$, define the cross-polar map
\begin{equation}\label{eq:cross-polar-map}
 \Phi_{b,L}:V\longrightarrow L^*,
 \qquad \Phi_{b,L}(x)(y)=b\cdot\beta(x,y),
\end{equation}
and let $e_L(b)=\rank\Phi_{b,L}$.  Put
\begin{equation}\label{eq:theta-def}
 e_* = \min_{b\ne0}e_L(b),
 \qquad
 \Theta_L(g)=\sum_{b\ne0}q^{-e_L(b)}.
\end{equation}
For $y\in L\setminus\{0\}$, also define the directional derivative
\begin{equation}\label{eq:directional-map}
 T_y:V\longrightarrow\F_q^K,
 \qquad T_y(x)=\beta(x,y),
\end{equation}
write $r(y)=\rank T_y$, and let
$\kappa(y)$ be one when $-(g(y)-g(0))\in\operatorname{im}T_y$ and zero
otherwise.  The indicator $\kappa(y)$ records whether two points separated by
$y$ can collide at all; the rank $r(y)$ then determines how many such pairs
exist.

\begin{theorem}[Root collisions and the rank-spectrum bound]
\label{thm:coset-success}
Choose an affine coset $C$ uniformly from $V/L$ and choose
$z\in\F_q^K$ uniformly.  Let
\[
 Z=|\{x\in C:g(x)=z\}|,
 \qquad \mu=q^{h-K},
\]
and define
\begin{equation}\label{eq:eta-def}
 \eta_0=\sum_{y\in L\setminus\{0\}}\kappa(y)q^{-r(y)},
 \qquad
 \bar\eta=\sum_{y\in L\setminus\{0\}}q^{-r(y)}.
\end{equation}
Then
\begin{align}
 \mathbb E Z&=\mu,\label{eq:EZ}\\
 \mathbb E[Z(Z-1)]&=\mu\eta_0,\label{eq:factorial-moment}\\
 \eta_0&\le\bar\eta
 =\mu\bigl(1+\Theta_L(g)\bigr)-1.\label{eq:eta-spectrum}
\end{align}
Consequently,
\begin{align}
 \operatorname{Var}(Z)&\le\mu^2\Theta_L(g),\label{eq:theta-var}\\
 \Pr[Z>0]&\ge\frac{\mu}{1+\eta_0}
 \ge\frac{1}{1+\Theta_L(g)}.\label{eq:coset-success}
\end{align}
If $h=K-\tau$ and $e_*\ge h-\Delta$, then
\begin{equation}\label{eq:defect-retry}
 \Theta_L(g)<q^{\tau+\Delta},
 \qquad
 \mathbb E[\textnormal{trials to a root-containing target--coset pair}]
 <1+q^{\tau+\Delta}.
\end{equation}
This last bound controls the search for a root-containing target--coset pair,
not uniqueness of the root.
Here $\Delta$ is the maximum loss from full cross-polar rank $h$.
\end{theorem}

\begin{proof}
See \Cref{app:proof-coset-success}.
\end{proof}

\paragraph{How to read the theorem.}
The first moment $\mathbb EZ=\mu$ is exactly the random-map value.  The second
factorial moment $\mathbb E[Z(Z-1)]$ counts ordered pairs of distinct roots and
therefore measures collisions.  The compatible collision mass $\eta_0$ is
exact; $\bar\eta$ drops the compatibility test and is a convenient upper
bound.  The sum $\Theta_L(g)$ describes the entire cross-polar rank spectrum,
not merely its minimum.  A few isolated rank defects may have little effect,
whereas low rank in many output directions can substantially increase retries.

\begin{center}
\begin{tabularx}{0.92\textwidth}{@{}>{\bfseries}p{0.27\textwidth}X@{}}
$\mu=q^{h-K}$ & expected number of roots in a random target--slice pair\\
$\eta_0$ & exact expected number of compatible collisions per planted root\\
$\bar\eta$ & rank-only upper bound on that collision mass\\
full cross-polar rank & for the selected positive deficits, unique-root trial
count essentially $q^{K-h}=q^\tau$
\end{tabularx}
\end{center}

\begin{corollary}[Unique-root trials and the planted distribution]
\label{cor:singleton-coset}
In the setting of \Cref{thm:coset-success},
\begin{align}
 \Pr[Z=1]&\ge\mu(1-\eta_0)
             \ge\mu(1-\bar\eta),\label{eq:singleton-coset}\\
 \Pr[Z>0]&\ge\mu\left(1-\frac{\bar\eta}{2}\right),
 \qquad (\bar\eta<2),\label{eq:bonferroni-success}\\
 \Pr[Z\ge2\mid Z>0]&\le\frac{\bar\eta}{2-\bar\eta},
 \qquad (\bar\eta<2).\label{eq:conditional-multiple}
\end{align}
In particular, if $\bar\eta<1$, the expected number of independent trials
until $Z=1$ is at most
\begin{equation}\label{eq:unique-retry-general}
 \frac{1}{\mu(1-\bar\eta)}.
\end{equation}

Let $\mathsf U$ be the uniform target--coset pair conditioned on $Z>0$, and
let $\mathsf P$ be the planted distribution obtained by choosing a uniform
coset, a uniform point in it, and the point's image as target.  If
$\bar\eta<1$, then
\begin{equation}\label{eq:planted-tv}
 \Pr_{\mathsf P}[Z\ge2]\le\eta_0\le\bar\eta,
 \qquad
 d_{\rm TV}(\mathsf U,\mathsf P)\le\bar\eta.
\end{equation}
\end{corollary}

\begin{proof}
See \Cref{app:proof-singleton-coset}.
\end{proof}

Total variation distance is the largest discrepancy that the two distributions
can assign to any event.  Thus, when $\bar\eta$ is tiny, an experiment that
plants a known root samples essentially the same target--coset distribution as
a uniformly chosen trial conditioned on being solvable.  This justifies using
planted instances to probe solver behavior; it does not remove the separate
finite-size-to-production extrapolation.

\begin{corollary}[Full cross-polar rank]\label{cor:full-cross-rank}
The following are equivalent:
\begin{enumerate}[label=(\roman*)]
\item $e_L(b)=h$ for every $b\ne0$;
\item $T_y$ is surjective for every $y\in L\setminus\{0\}$.
\end{enumerate}
Under these conditions,
\[
 \eta_0=\bar\eta=\mu-q^{-K}.
\]
If $h=K-\tau$ with $\tau\ge0$, the expected number of trials until a unique
base-field root is at most
\begin{equation}\label{eq:unique-retry}
 \mathsf{Retry}_{K,\tau}\coloneqq
 \frac{q^\tau}{1-q^{-\tau}+q^{-K}}.
\end{equation}
\end{corollary}

\begin{proof}
If some $T_y$ is not surjective, a nonzero $b$ annihilates its image, so the
transpose of $\Phi_{b,L}$ has the nonzero kernel vector $y$ and $e_L(b)<h$.  The
converse is the same argument in reverse.  Under either equivalent condition,
$r(y)=K$ and $\kappa(y)=1$ for every nonzero $y$, giving
$\eta_0=\bar\eta=(q^h-1)q^{-K}=\mu-q^{-K}$.  Substitute this value into
\eqref{eq:unique-retry-general}.
\end{proof}

For the positive deficits used in the concrete attacks, the expected
unique-root trial count under full cross-polar rank is essentially $q^\tau$:
the denominator in \eqref{eq:unique-retry} differs from one only by the tiny
terms $q^{-\tau}$ and $q^{-K}$.  In the running Level-I example, $h=40$, $K=50$,
and $\tau=10$, so the expected number of independent target--slice trials is
essentially $19^{10}=2^{42.48}$.  Conditional on a root-containing trial, the
probability of more than one base-field root is below $2^{-43.47}$.  This is
the precise justification for the retry factor used in the structured cost
model.

Equation~\eqref{eq:eta-spectrum} distinguishes sparse rank defects from a
pervasive loss.  If at most $T$ nonzero coefficient vectors $b$ are defective
and each has defect $h-e_L(b)\le\Delta$, then
\begin{equation}\label{eq:defect-spectrum}
 \bar\eta\le
 \mu-q^{-K}+Tq^{-K}(q^\Delta-1).
\end{equation}
If defects are counted projectively, $T$ in this formula must include all
$q-1$ nonzero scalar multiples of each projective direction.  Under full
cross-polar rank and $q=19$, \eqref{eq:conditional-multiple} is below
$2^{-43.47}$ for $\tau=10$ and below $2^{-60.47}$ for $\tau=14$; the
corresponding total-variation bounds are below $2^{-42.47}$ and $2^{-59.47}$.

\begin{corollary}[Full-space target availability]\label{cor:full-space}
Let $g:\F_q^N\to\F_q^K$ have degree at most two.  Let $r_*$ be the
minimum polar rank of a nonzero output combination, and put
$\alpha=1-q^{-K}$.  For a uniform target,
\begin{equation}\label{eq:full-space-success}
 \Pr[g^{-1}(z)\ne\varnothing]
 \ge \frac{1}{1+\alpha q^{K-r_*}}.
\end{equation}
Hence the expected number of target trials is at most
$1+\alpha q^{K-r_*}$.  If $r_*\ge2K$, every target has a root.
\end{corollary}

\begin{proof}
See \Cref{app:proof-full-space}.
\end{proof}

For a structured $\ell=4$ lift, an invertible output transformation separates
the $M-K$ affine coordinates from the remaining $K$ coordinates of a uniform
target.  Conditioning on the affine coordinates fixes a base point; the other
coordinates remain uniform, and translation does not change the polar maps.
We take $V=\ker C_{R,\mathbf v}$ and let $L$ be the selected stable slice, of dimension
$h=\ell s=K-\tau$.  Full cross-polar rank gives the unique-root retry bound
$\mathsf{Retry}_{K,\tau}$ in \eqref{eq:unique-retry}.  A worst-case
rank-defect bound gives the budget for reaching a root-containing target--coset
pair in \eqref{eq:defect-retry}; unique recovery instead
depends on the aggregate quantity $\bar\eta$ and the Macaulay kernel model.

For the fixed-skew $\ell=2$ lift, take $L=V$ to be the complete residual space.
By \Cref{prop:fixed-skew}, every root is accepted.  The expected target-search
factor is at most $1+\alpha q^{K-r_*}$, and $r_*\ge2K$ makes every target
solvable.  Under the AXN conversion, the arithmetic budgets remain below the
nominal NIST references for global ranks at least 46, 68, and 89 at Levels I,
III, and V.  The sufficient all-target thresholds are 96, 144, and 192.  These
are global conditions and concern the converted operation model, not a
complete solver circuit; the sampled evidence in \Cref{sec:evidence} is not a
certificate.  The Level-V filter attack instead uses
\Cref{prop:l2-rejection}.

\subsection{Sampling exactly uniform targets with the official conversion}

The probability theorems above assume a uniform target in $\F_q^K$.
Modulo reduction of binary hash chunks can introduce a small bias, so we do
not simply assume uniformity.  The verifier need not be changed to remove the
bias.  The attacker repeatedly chooses a
message--salt input, runs the official hash-to-field conversion, and keeps the
input only when each underlying binary chunk lies in a range that makes the
resulting base-$19$ digits exactly uniform.

Concretely, the official $q=19$ conversion reads at most fifteen base-19 digits
from each 64-bit chunk.  If a uniform $b$-bit chunk $X$ is to supply $a$
digits, the attacker accepts that message--salt candidate only if
\begin{equation}\label{eq:uniform-chunk}
 X<\left\lfloor\frac{2^b}{q^a}\right\rfloor q^a.
\end{equation}
Conditioned on this event, $X\bmod q^a$ is uniform, so the extracted
digits are jointly uniform.  Applying \eqref{eq:uniform-chunk} independently
to every chunk makes the retained official target exactly uniform.  The
invertible output-coordinate change used in \Cref{sec:affine-columns}
preserves uniformity.  Conditioning on the values of the $M-K$ affine-output coordinates---or on
their passing values in a filtered lift---then leaves the remaining $K$
residual target coordinates exactly uniform, which is the distribution
assumed in the root theorems.  This is
ordinary rejection sampling performed by the forger over message--salt
candidates; the verifier continues to apply the unmodified official
conversion.  Across all Version~2.3 lengths the acceptance probability is at
least $0.152462$, or fewer than 6.6 candidates on average.

The Level-V $\ell=2$ filter needs about $19^{32}>2^{128}$ raw targets before
the $M-K=32$ consistency conditions pass, so the 128-bit salt alone does not
supply enough distinct inputs.  An existential forger may vary a chosen-message
prefix as well as the salt.  Distinct prefix--salt pairs are distinct random-oracle
inputs and therefore give independent targets in the model.  The
artifact's deliberately loose target-generation envelope is below $2^{32}$
gates per target retained by the unbiased hash-chunk sampler (the AXN charge
is no larger than the NAND envelope used there).  Even after multiplication by the $19^{32}$ filtering factor,
this term is more than 70 exponent bits below the cost-driving MQ solve.

\section{Concrete Costs}\label{sec:complexity}

The algebraic reduction determines the residual MQ systems exactly.  Turning
those systems into numerical security estimates requires a separate cost
model.  We report that model in layers so that a reader can see which numbers
are structural, which come from an MQ solver estimate, and which are merely a
change of arithmetic units.

\begin{table}[H]
\centering
\caption{How to interpret the cost accounting.}
\label{tab:cost-layers}
\small
\renewcommand{\arraystretch}{1.12}
\begin{tabularx}{0.96\textwidth}{@{}>{\bfseries}p{0.23\textwidth}X@{}}
\toprule
Layer & What it means\\
\midrule
Residual system & Exact numbers of equations, variables, blocks, and public
rank conditions after the reduction.\\
Solver model & Published generic-MQ or special-XL formula for the dominant
finite-field work at those dimensions.\\
AXN conversion & Replacement of the solver formula's coarse field-operation
factor by an explicit two-input AND/XOR/XNOR circuit for one signed
multiply-accumulate.\\
Not included & A complete implementation of sparse matrix construction,
pivoting and control, memory addressing and traffic, communication, or an
optimized end-to-end attack netlist.\\
\bottomrule
\end{tabularx}
\end{table}

A cost exponent $c$ means modeled work proportional to $2^c$ in the stated
unit.  A difference of ten exponent bits is a factor of about $2^{10}$, not ten
percent.  Because the solver formulas count dominant arithmetic rather than a
complete machine, we call the difference from a NIST reference an
\emph{arithmetic overhead budget}.

\subsection{Cost units and circuit conversion}\label{sec:units}

The cited attacks first count finite-field operations and then apply coarse
bit-cost factors~\cite{Beullens2024SNOVA,Cabarcas2025SNOVA}.  NIST lists
$2^{143}$, $2^{207}$, and $2^{272}$ as nominal classical-gate references, while
also requiring category comparisons across every metric it considers relevant
to practical security~\cite{NISTPQCCriteria}.  We retain each attack's
field-operation count and replace only its dominant field factor by explicit
circuits.  The resulting figures sharpen one metric; they do not by themselves
establish a category break.  Our primary AXN basis assigns unit cost to two-input
AND, XOR, and XNOR gates, with free constants, wires, fan-out, and structural
sharing.  This is the gate family used by the concrete AES-128 circuit on
NIST's circuit list~\cite{NISTCircuitComplexity}.  A two-input-NAND conversion
provides a separate basis-sensitivity check.

For generic MQ, the published factor is
\begin{equation}\label{eq:published-f19-factor}
 b_1=(\log_2 19)^2+\log_2 19=22.29.
\end{equation}
Using canonical five-bit representatives and exact Barrett reduction, the
validated circuits for one signed multiply-accumulate, $c+ab$ or $c-ab$, over
$\F_{19}$ use 402 and 441 AXN gates, respectively; their NAND translations use 781 and 835 gates.  To avoid
assuming a particular sign schedule in elimination, we charge
\begin{equation}\label{eq:axn-f19-factor}
 g_1^{\mathrm{AXN}}=441,
 \qquad g_1^{\mathrm{NAND}}=835.
\end{equation}
Both signs are checked exhaustively on all $19^3$ valid input triples by the
generator and by an independent scalar evaluator.

For special XL, we work in the isomorphic sparse basis
$\F_{19}[t]/(t^4-t-1)$.  In this basis, the element
$13+13t+15t^2+2t^3$ is a root of the official defining polynomial; evaluating
its powers supplies the explicit basis change to $\F_{19}[S]$.  Two-level
Karatsuba uses nine base-field multiplications, and $t^4=t+1$ makes reduction
additive.  The validated circuits for $c\mathbin{\pm}ab$ over $\F_{19^4}$ use
6,077 and 6,081 AXN gates, or 11,086 and 11,102 NAND gates.  We therefore
charge
\begin{equation}\label{eq:axn-f194-factor}
 g_4^{\mathrm{AXN}}=6{,}081,
 \qquad g_4^{\mathrm{NAND}}=11{,}102.
\end{equation}
The artifact checks the complete scalar basis tensor, 4,096 random extension
products, the basis isomorphism, and both translated netlists.  Charging a
signed multiply-accumulate for every multiplication in the special-XL formula
is conservative.

\paragraph{Same-basis AES calibration.}
The release builds variable-key AES-128/192/256 one-pair key-filter circuits
with the NIST 113-gate S-box, the complete key schedule, encryption, and a
final ciphertext comparator.  To identify the actual key rather than merely a
key that matches one pair, surviving candidates are checked lazily on enough
independent plaintext--ciphertext pairs to leave at most $2^{-128}$ expected
wrong survivors in the ideal-cipher model: a total of two pairs for AES-128 and
a total of three for AES-192 and AES-256.  The additional average work is
negligible because only
survivors reach those checks.  Every first-pair circuit is validated on the
FIPS test vector and eight additional keys against an independent AES
implementation.

\begin{table}[tbp]
\centering
\caption{Constructive AES calibration circuits.  Full charges the entire key
space for the first-pair filter; expected search is one bit lower, and the
lazy confirmation overhead is negligible under the stated ideal-cipher model.
These circuits are not lower bounds on optimal AES implementations and do not
replace NIST's nominal references.}
\label{tab:aes-calibration}
\scriptsize
\setlength{\tabcolsep}{3.5pt}
\begin{tabular}{@{}c r r r r r r@{}}
\toprule
Level & pairs & NIST nominal & AXN gates/filter & AXN full & NAND gates/filter & NAND full\\
\midrule
I   & 2 & 143 & 30,081 & 142.88 & 106,616 & 144.70\\
III & 3 & 207 & 34,227 & 207.06 & 121,735 & 208.89\\
V   & 3 & 272 & 41,543 & 271.34 & 147,772 & 273.17\\
\bottomrule
\end{tabular}
\end{table}

The NIST-listed 28,600-gate AES-128 circuit has full-key-space exponent
$128+\log_2(28{,}600)=142.80$, numerically consistent with the nominal value
143.  Across our nine rows, the AXN arithmetic estimates are
\AXNFullAESGapRange{} bits below the constructive full-space AES circuits and
\AXNNominalHeadroomRange{} bits below the nominal NIST references.  The NAND
cross-check gives attack ranges \LoneNANDRange, \LthreeNANDRange, and
\LfiveNANDRange.

\paragraph{Scope of the conversion.}
The AXN and NAND figures price the field-arithmetic schedule counted by the
published estimators.  They do not synthesize pivot selection, solver control,
address generation, memory traffic, data movement, or a complete
sparse-linear-algebra implementation.  They are therefore constructive
conversions of inherited arithmetic models, not complete attack netlists.
Public preprocessing and exact-uniform target generation are bounded
separately.  The dominant unresolved inputs are the global production rank
spectra and the accuracy of the production-scale solving models.

\subsection{Structured special XL for \texorpdfstring{$\ell=4$}{l=4}}\label{sec:special-xl}

The exact reduction supplies an $A$-linear slice of $A$-dimension $s$,
hence base-field dimension $4s$ when $\ell=4$.  After extension to the
splitting field, this becomes four blocks of $s$ variables.  For each unordered pair of blocks there are $m_1$ independent equations
of degree two: diagonal equations use one block twice, and off-diagonal
equations use two blocks once each.  Special XL records a separate degree for
each block and selects a multidegree
at which the associated Macaulay matrix is predicted to have the needed rank.
A reader need not manipulate the series below: here it is only a compact device
for choosing the matrix degree and counting the resulting monomial columns.

After adding one homogenizing variable to each block, the Hilbert series used
by the published special-XL model is
\begin{equation}\label{eq:hilbert-series}
 H_{m_1,s}(\mathbf t)=
 \frac{\prod_{a=1}^{4}(1-t_a^2)^{m_1}
       \prod_{1\le a<b\le4}(1-t_at_b)^{m_1}}
      {\prod_{a=1}^{4}(1-t_a)^{s+1}}.
\end{equation}
For $\mathbf d=(d_1,\ldots,d_4)$, put
\begin{equation}\label{eq:macaulay-size}
 M_s(\mathbf d)=\prod_{a=1}^{4}\binom{s+d_a}{d_a}.
\end{equation}
The search chooses a multidegree whose signed Hilbert coefficient is
negative.  This sign change is a rank-prediction trigger in the cited
special-XL model: a random Macaulay matrix is predicted to have full column
rank, while a system with a planted solution is predicted to have a
one-dimensional right kernel
\cite{Cabarcas2025SNOVA}.  The modeled linear-algebra work for one target--slice trial is
$3s^2\bigl(M_s(\mathbf d)\bigr)^2$ extension-field multiplications: the square is the
dense linear-algebra term in the published model, and $3s^2$ is its stated
prefactor.

Cabarcas et al.\ handle an underdetermined system by specializing $k$
coordinates, where $4\mid k$, and multiplying the work by $q^k$.  Our attack
uses a different retry loop: it fixes a $4s$-dimensional slice and samples the
target and affine coset.  Under full cross-polar rank,
\Cref{cor:full-cross-rank} bounds the expected trials needed for a unique
base-field root by $\mathsf{Retry}_{K,\tau}$.  The modeled work is therefore
\begin{equation}\label{eq:special-xl-cost}
 \mathsf{Retry}_{K,\tau}\,3s^2M_s(\mathbf d)^2,
 \qquad \tau=K-4s,
\end{equation}
extension-field multiplications.  Replacing the older $1+q^\tau$ factor by
$\mathsf{Retry}_{K,\tau}$ changes each displayed logarithm by less than $10^{-25}$ bits.
If every cross-polar map has rank defect at most $\Delta$, then
\eqref{eq:defect-retry} bounds the expected target--coset trials---and hence
the associated Macaulay linear-algebra budget under the same per-trial
model---by $1+q^{\tau+\Delta}$.  It does not by itself imply unique
recovery.  That step depends on $\bar\eta$ and
on the Macaulay kernel behavior described below.

The computation takes place over $\F_{19^4}$.  The published convention assigns
each extension-field multiplication
\begin{equation}\label{eq:bitop-conversion}
 b_4=2\bigl(\log_2 19^4\bigr)^2+\log_2 19^4=594.43
\end{equation}
bit operations.  The AXN arithmetic conversion replaces $b_4$ by the conservative
signed multiply-accumulate bound $g_4^{\mathrm{AXN}}$ from
\eqref{eq:axn-f194-factor}.

\begin{lemma}[Recovery from a one-dimensional Macaulay kernel]\label{lem:macaulay-recovery}
Fix a target $t$, a selected affine slice $a(t)+y+L$, and an
$\F_q$-linear parameterization $\iota:\F_q^h\to L$.  Let $\lambda_*\in\F_q^h$ be slice coordinates for a
root
\[
  u_*=a(t)+y+\iota(\lambda_*),
\]
and let $x_*\in\Omega^h$ be the image of $\lambda_*$ under the eigenblock
change of variables.  Let $\mathcal M$ be the corresponding special-XL
Macaulay matrix over the splitting field $\Omega$, dehomogenized by setting
every block homogenizer to one.  Suppose its columns include the constant and
degree-one monomials.  If $\ker_\Omega\mathcal M$ is one-dimensional, row
reduction recovers $x_*$, $\lambda_*$, the common column $u_*$, and hence the
complete signature.
\end{lemma}

\begin{proof}
The dehomogenized monomial-evaluation vector $\nu(x_*)$ lies in
$\ker_\Omega\mathcal M$ and has constant coordinate one.  Every nonzero
kernel vector is therefore a scalar multiple of $\nu(x_*)$.  Normalizing its
constant coordinate to one recovers $x_*$ from the degree-one coordinates.  The
inverse eigenblock transformation recovers $\lambda_*\in\F_q^h$; the affine
embedding gives $u_*=a(t)+y+\iota(\lambda_*)$; and the column relation
\eqref{eq:relation} reconstructs every signature column.
\end{proof}

\paragraph{Las Vegas wrapper.}
For each target--coset trial, compute the Macaulay kernel and continue unless
it is one-dimensional with nonzero constant coordinate.  Normalize that
coordinate, recover the degree-one eigenblock coordinates, and invert the
eigenblock transformation.  Continue again unless the point lies in $\F_q$
and the public verifier accepts it.  The procedure therefore never returns an
invalid signature; the solving model predicts only how often a trial reaches
the recoverable case with a one-dimensional right kernel.


The target--coset theorem counts $\F_q$-rational roots, so a unique-root
trial contains the point required by the lemma.  \Cref{cor:singleton-coset}
bounds multiple-root trials and the distance from the planted-root
distribution using the full cross-polar rank spectrum.  Under full
cross-polar rank, the conditional multiple-root probability is below
$2^{-43.47}$ for $\tau=10$ and below $2^{-60.47}$ for $\tau=14$.

The remaining assumption that the production Macaulay right kernel is
one-dimensional concerns extra algebraic kernel vectors, including those
caused by extension-field roots or by the Macaulay rows failing to span all
expected relations; it is not an uncharged descent retry.  Every recovered point is checked by the original
verifier.  The small-profile tests in \Cref{sec:evidence} support, but do not
prove, this production solving model.

\paragraph{Precisely what the displayed special-XL costs assume.}
The exponents in \Cref{tab:xl-optima} use (i) global full cross-polar rank on
the selected slice, so that $\mathsf{Retry}_{K,\tau}$ bounds unique-root trials, and (ii) the
production Macaulay behavior just described: a planted-solvable trial has a
one-dimensional right kernel.  A global rank-defect bound alone controls only
retries to a root-containing target--coset pair through
\eqref{eq:defect-retry}; it does not imply
unique recovery or a one-dimensional right kernel.  This
separation is important because the first condition concerns an entire
family of nonzero output directions (modulo scalar multiples), whereas the second concerns production-scale algebraic
solving behavior.

\begin{table}[tbp]
\centering
\caption{Cost-driving special-XL parameters for the six $\ell=4$ rows.
``Model'' uses $b_4$; ``AXN arithmetic'' replaces it by the validated
signed multiply-accumulate circuit $g_4^{\mathrm{AXN}}$.}
\label{tab:xl-optima}
\scriptsize
\setlength{\tabcolsep}{4pt}
\begin{tabular}{@{}c c c c c r r@{}}
\toprule
$m_1$ & available $\dim_A H$ & $K$ & $(s,\tau)$ & $\mathbf d$ & model & AXN arithmetic\\
\midrule
5 & $12$--$13$ & 50 & $(10,10)$ & $(3,3,3,4)$ & 128.82 & 132.17\\
7 & $15$--$19$ & 70 & $(15,10)$ & $(4,4,6,6)$ & 171.69 & 175.04\\
9 & $22$--$23$ & 90 & $(19,14)$ & $(5,6,6,6)$ & 214.12 & 217.48\\
\bottomrule
\end{tabular}
\end{table}

The exact search checks every admissible $s$ and all 42,356 sorted
multidegrees whose modeled work could match or beat the incumbent, including
highly unbalanced degrees.  It recovers the three optima in
\Cref{tab:xl-optima}, uniquely up to block permutation.  The conservative
target-generation charge does not change the displayed hundredths.

For the search for a root-containing target--coset pair and the associated
Macaulay linear algebra, the AXN budget retains
a positive nominal arithmetic budget for $\Delta\le2,7,12$ at Levels I, III,
and V.
These thresholds neither price a complete solver circuit nor certify unique
recovery.  The latter additionally requires $\bar\eta<1$ in
\eqref{eq:eta-def} and the modeled production Macaulay matrix to have a
one-dimensional right kernel.

\subsection{The \texorpdfstring{$\ell=2$}{l=2} attacks}

For Levels I and III, use the fixed-skew lift of \Cref{prop:fixed-skew}.  The
public preflight gives quotient rank $K$, affine rank $M-K$, and residual
systems of dimensions $48/112$ and $72/168$.  CFXL is a hybrid strategy
that specializes some variables and then applies XL-style linear algebra;
Hashimoto refines this optimization for underdetermined MQ.\@  Minimizing these
variants in Beullens's printed Appendix-A convention gives the per-solve
estimates
\begin{equation}\label{eq:fixed-skew-costs}
 2^{129.44}\quad\text{and}\quad2^{183.44}
\end{equation}
modeled bit operations~\cite{Beullens2024SNOVA,Hashimoto2023}.  Replacing
$b_1$ by $g_1^{\mathrm{AXN}}$ gives arithmetic exponents $133.75$ and $187.74$.
Every recovered root is accepted.  The headline per-solve estimates incur
no target-retry factor when the global all-target conditions $r_*\ge96$ and
$r_*\ge144$ hold.  For a smaller
global rank, the expected-cost exponent increases by at most
$\log_2(1+(1-19^{-K})19^{K-r_*})$.  The same fixed-skew construction at
Level~V gives a $96/224$ fallback with model exponent $237.84$ and AXN exponent
$242.15$; its all-target threshold is $r_*\ge192$.

The optimizer includes Hashimoto's published boundary value $a=1$.  For
$\operatorname{MQ}(19,1,1)$, the semi-regular Hilbert-series rule used by CFXL
never produces the stopping degree at which its linearization would be costed.
The artifact therefore solves that elementary system by 19 direct evaluations
and deliberately overcharges each evaluation.  This
lowers six generic quotient baselines and the Level-V fixed-skew fallback; it
does not change the selected Level-I or Level-III fixed-skew costs.

The Level-V headline instead uses $R_0=I,R_1=0$ and zero offsets.  The quotient
has rank $K$, target filtering enforces the $M-K$ public consistency
conditions, and the residual system has 96 equations in all 256 variables.
Target search occurs before equation solving, so its cost adds to---rather than
multiplying---the cost of a solve.  The generic-MQ solve has AXN arithmetic
exponent $239.50$, while exact-uniform target filtering is smaller by more than
70 exponent bits.  The displayed exponent incurs no target-retry factor under
the accepted-root condition \eqref{eq:l2-rejection}; for this row, $r_*\ge240$
suffices.  At lower rank, \eqref{eq:l2-accepted-target} gives the additional expected retry factor.
The fixed-skew fallback is slower but makes every recovered root acceptable by
construction.

\begin{table}[tbp]
\centering
\caption{$\ell=2$ systems and global polar-rank thresholds.  Model is the
pinned published operation model; AXN uses $g_1^{\mathrm{AXN}}=441$.  ``For
budget'' is the smallest rank for which expected AXN arithmetic remains below
the nominal NIST reference.  ``All-target root'' is sufficient to guarantee
that every target has an accepted root; it does not assert that one particular
solver invocation succeeds.}
\label{tab:l2-costs}
\scriptsize
\setlength{\tabcolsep}{3.25pt}
\begin{tabular}{@{}l l r r r r r@{}}
\toprule
Parameters & mode & eqs./vars. & model & AXN & \shortstack{for\\budget} & \shortstack{all-target\\root}\\
\midrule
$(48,16,2,2)$ & fixed-skew & $48/112$ & 129.44 & 133.75 & 46  & 96\\
$(72,24,2,2)$ & fixed-skew & $72/168$ & 183.44 & 187.74 & 68  & 144\\
$(96,32,2,2)$ & filter     & $96/256$ & 235.19 & 239.50 & 177 & 240\\
$(96,32,2,2)$ & fixed-skew & $96/224$ & 237.84 & 242.15 & 89  & 192\\
\bottomrule
\end{tabular}
\end{table}

For the filter row, the all-target-root threshold includes rejection through
\eqref{eq:l2-rejection}; for fixed skew, it is the simpler condition
$r_*\ge2K$ from \Cref{cor:full-space}.  The quotient matrix depends only on the
fixed public outer map and the relation, not on a generated key.  Exhaustive
enumeration of all $19^2=361$ relations $R_1=aI+bS$ proves quotient rank $K$
for every relation on all three parameter shapes.  Exhaustive optimization of
the natural two-block special-XL model gives bit-operation exponents $132.31$,
$184.27$, and $236.36$, all worse than the selected generic-MQ modes.

\section{Implementation and Evidence}\label{sec:evidence}

The artifact implements the public key generator, verifier-side feature maps,
symmetry quotient, affine reduction, structured lift, rank tests, and cost
searches over $\F_{19}$.  It does \emph{not} report a production-size
end-to-end forgery, a production-size Macaulay solve, or a wall-clock attack
timing.  Its full-size experiments concern public preprocessing and rank; its
solver experiments use smaller tractable profiles.  The checks therefore have
different logical force.  The table
below is the intended reading guide for every experimental statement in this
section.

\begin{table}[H]
\centering
\caption{Evidence hierarchy.  ``Exact'' here means exact for the concrete
object checked, not an exhaustive theorem about all production directions.}
\label{tab:evidence-hierarchy}
\small
\renewcommand{\arraystretch}{1.12}
\begin{tabularx}{0.96\textwidth}{@{}>{\raggedright\arraybackslash\bfseries}p{0.24\textwidth}>{\raggedright\arraybackslash}p{0.30\textwidth}>{\raggedright\arraybackslash}X@{}}
\toprule
Check & What is computed & What it does---and does not---establish\\
\midrule
Public preflight & Exact ranks and rejection conditions for the chosen key,
relation, offsets, and slice & The exact reduction is valid for that concrete
attack instance.\\
Rank campaign & Exact ranks for many selected nonzero output or direction
vectors & Finite evidence about the unenumerated global minimum rank; not a
certificate for all directions.\\
Macaulay experiment & Exact kernels at smaller tractable dimensions & Evidence
for the production solver model; not a production timing or rank proof.\\
Circuit validation & Exhaustive or independent checking of the released field
primitive & Exact cost and functionality of that primitive, but not of a
complete solver.\\
\bottomrule
\end{tabularx}
\end{table}

A ``projective direction'' below means a nonzero coefficient vector modulo
multiplication by a nonzero field scalar; all scalar multiples define the same
polar rank.  Sampling many distinct directions can expose low-rank structure,
but only exhaustive enumeration or a proof can certify the global minimum
required by the headline all-target-root conditions.

\subsection{Public preflight and verifier regression}

For a structured $\ell=4$ instance, the preflight checks
\[
 \rank Q_R=K,\qquad \rank C_{R,\mathbf v}=M-K,\qquad
 \rank\mathcal O_{R,\mathbf v}=m_1\ell^2,
\]
together with $\mathcal F_{R,\mathbf v}H_{R,\mathbf v}=0$, the predicted $A$-kernel
dimension, $A$-invariance, the restricted quadratic rank, and rejection
compatibility.  All 121 regenerated full-size instances tested in the artifact pass,
including the official Level-I known-answer key.  The first deterministic
offset candidate passes for 114 instances; the second passes for the remaining
seven.

For fixed skew, the preflight checks the relation and rejection identity and
the quotient, affine, and augmented ranks.  All 60 full-size keys tested in the artifact pass.  The artifact also reproduces the official Level-I public key byte for byte,
verifies its known-answer signature against the official hash target, and
checks the fixed output-order permutation described in
\Cref{sec:affine-columns}.  Eighteen additional regressions cover all nine
Version~2.3 shapes.  An independent Gauss--Jordan implementation agrees with
the rank kernel on 2,036 matrices, including 36 production-sized matrices.

\subsection{Rank evidence}

\paragraph{Structured slices.}
Every selected slice is tested on all coefficient-basis directions (one
nonzero output coefficient at a time) and 1,024 deterministic dense directions
(many nonzero coefficients).  Selected slices also receive complete normalized
weight-two tests (two nonzero coefficients, modulo overall scaling) or
additional dense tests.  All 102,796 tested
cross-polar maps have maximal rank $h=4s$; the per-shape counts are in
\Cref{tab:cross-polar}.  The dual campaign evaluates 21,992 directional maps
$T_y$ on the same 13 slices, and every one is surjective of rank $K$.
All 124,788 coefficient and direction vectors used in these two campaigns are
projectively distinct within their respective slice records.

\begin{table}[tbp]
\centering
\caption{Exactly computed cross-polar ranks for the tested structured $\ell=4$ slice--direction pairs.}
\label{tab:cross-polar}
\small
\setlength{\tabcolsep}{5.2pt}
\begin{tabular}{@{}l r r r r@{}}
\toprule
Parameters & slices & $h=4s$ & rank tests & sampled minimum\\
\midrule
$(28,5,4,4)$  & 3 & 40 & 26,296 & 40\\
$(28,4,4,5)$  & 2 & 40 & 24,198 & 40\\
$(40,7,4,4)$  & 2 & 60 & 45,658 & 60\\
$(38,5,4,5)$  & 2 & 60 & 2,188  & 60\\
$(50,9,4,4)$  & 2 & 76 & 2,228  & 76\\
$(52,6,4,6)$  & 2 & 76 & 2,228  & 76\\
\midrule
Total          & 13 & -- & 102,796 & maximal\\
\bottomrule
\end{tabular}
\end{table}

\paragraph{Fixed-skew systems.}
A broad campaign tests every coefficient-basis direction together with
combinations having two, three, or many nonzero output coefficients.  A second campaign targets
the 21 exceptional projective directions that make the local operator on
$\Mat_2(\F_{19})$ rank two rather than four.  It tests all 381 local
projective directions in one block, every exceptional direction in every
remaining block, and deterministic two-block exceptional combinations.
Together the campaigns perform 3,888 exact rank evaluations; their sampled
minima and the rigorous thresholds used by the cost analysis appear in
\Cref{tab:fixed-skew-evidence}.

\begin{table}[tbp]
\centering
\caption{Full-size polar-rank evidence for the fixed-skew $\ell=2$ systems.
``For budget'' is the minimum global polar rank that preserves positive AXN
arithmetic budget in \Cref{tab:l2-costs}; ``all-target root'' guarantees root
existence for every target, not success of a particular solver run.}
\label{tab:fixed-skew-evidence}
\small
\setlength{\tabcolsep}{4.5pt}
\begin{tabular}{@{}l r r r r r@{}}
\toprule
Parameters & preflights & rank tests & sampled min. & for budget & all-target root\\
\midrule
$(48,16,2,2)$ & 20 & 1,104 & 110 & 46 & 96\\
$(72,24,2,2)$ & 20 & 1,296 & 167 & 68 & 144\\
$(96,32,2,2)$ & 20 & 1,488 & 223 & 89 & 192\\
\midrule
Total & 60 & 3,888 & -- & -- & --\\
\bottomrule
\end{tabular}
\end{table}

\paragraph{Zero-offset filter.}
The broad and exceptional-direction campaigns perform 29,424 exact rank
evaluations on two keys per shape.  Their sampled minima are 127, 190, and 254
in dimensions 128, 192, and 256.  For the Level-V filter, the sampled minimum 254 is above the sufficient
accepted-root threshold 240.  If the global minimum is at least 254, then
\eqref{eq:l2-accepted-fraction} gives an accepted fraction of at least
$0.999999999$ in every target fiber.  The experiment supports that hypothesis
but does not certify the global minimum.

\subsection{Solver and probability checks}

The special-XL search examines all 42,356 admissible multidegrees at modeled
work no greater than the incumbent and recovers the three optima in
\Cref{tab:xl-optima}.  At tractable sizes, the artifact builds 128 four-block
Macaulay matrices: 64 random instances have full column rank and 64
planted-solvable instances have a one-dimensional right kernel, with the
planted evaluation vector in that kernel.  A 48-matrix two-block control
contains one planted case with a two-dimensional right kernel, so the
corresponding $\ell=2$ special-XL estimates are not used.  These are generic small-profile
experiments, not production SNOVA timings.

Exhaustive checks over $\F_3$ and $\F_5$ verify the exact first and factorial
second moments, the compatible-collision formula, the rank-spectrum upper bound,
and the singleton and planted-distribution inequalities on 110 complete
quadratic maps.  Seventeen maps have nonmaximal cross-polar rank; in twelve,
compatibility makes $\eta_0$ strictly smaller than $\bar\eta$.  Finally, on six actual cost-driving $\ell=2$ cores, the homogeneous
quadratic outputs are linearly independent, the common polar radical---the set
of directions annihilated by every polar form---is zero, and all 384 sampled
homogeneous Jacobians (matrices of first derivatives of the homogeneous
quadratic parts) attain the maximum rank permitted by their dimensions.  These
are genericity diagnostics, not solving-degree measurements.

\subsection{Scope of the estimates}

Once a public preflight passes, the quotient, affine reduction, stable kernel,
fixed-skew identity, and rejection handling are exact.  The remaining
production-scale inputs are the global polar spectra and the solving behavior.
The sampled ranks are broad but finite evidence and do not enumerate the full
families of nonzero output directions (modulo scalar multiples).  Likewise, the displayed exponents use the published special-XL and
generic-MQ models rather than full-size Gr\"obner or XL runs; no
production-size end-to-end solve is claimed.  The artifact
checks every dimension, Hilbert coefficient, public rank, and cost supplied to
those models and verifies every recovered candidate with the original public
verifier.

\paragraph{What would turn the estimates into end-to-end certificates.}
The rank assumptions would require either an algebraic proof of the global
minimum over every nonzero output direction or a complete enumeration enabled
by additional structure.  The solver assumptions would require
production-size Macaulay or Gr\"obner runs, or a validated implementation
whose observed ranks and costs replace the semi-regular models.  A complete
gate comparison would additionally need to price control flow, pivoting,
indexing, memory, and data movement.

\section{Countermeasures}\label{sec:countermeasures}

The design lesson is that security must be based on distinct quadratic
polynomials, not on the number of ordered labels used to name them.  A security
analysis for any symmetric redesign should compute the actual quotient rank
$\rho_{\rm quad}$ and use the unconditional cap
\begin{equation}\label{eq:countermeasure-K}
 \rho_{\rm quad}\le K_d
 =\min\left\{M,m_1\binom{d+1}{2}\right\},
 \qquad d=\dim_{\F_q}\F_q[S].
\end{equation}
Equality with $K_d$ is a public rank condition, not a consequence of symmetry
alone.  Outer randomization does not restore the alternating directions removed
before the outer map.

The direct repairs are to increase parameters after accounting for
\eqref{eq:countermeasure-K},
remove public-matrix symmetry, or redesign the left and right power actions so
that exchanging their labels changes the quadratic polynomial.  The
symmetric-block rejection rule is not sufficient: the fixed-skew $\ell=2$
lift makes every block nonsymmetric by construction, and the structured
$\ell=4$ preflight can certify the same property block by block.  Returning
unchanged to the submitted $q=16$ design would instead restore the known
characteristic-two attack surface.

A revised security analysis should publish the quotient rank, affine or filter
rank, stable-kernel rank when applicable, the polar or cross-polar condition
used for target availability, and solver logs for the actual cost-driving
system.

\section{Conclusion}\label{sec:conclusion}

The central vulnerability is an exact duplicate-feature identity.  \SNOVA{}
labels a quadratic power-pair feature by an ordered pair $(\xi,\eta)$, but
symmetry of the public base matrices makes that polynomial identical to the
one labeled $(\eta,\xi)$.  The public $q=19$ verifier therefore has at most
$m_1\binom{\ell+1}{2}$ distinct power-pair features where the ordered
representation has $m_1\ell^2$.  Because the identification occurs before the
public outer map, no later randomization can restore the missing directions.

The rest of the attack turns that identity into complete forgery systems rather
than merely a distinguisher or a reduced subset of verifier equations.  Exact
affine elimination retains every output.  A stable $A$-linear slice preserves
the four-block structure needed for the six $\ell=4$ special-XL systems.
Fixed-skew and target-filtered relations handle the three $\ell=2$ systems and
their block-symmetry rejection rule.  Exact moment formulas distinguish root
existence, unique-root retries, and the separate question of whether the
chosen Macaulay matrix exposes the root.  Every candidate is checked by the
unmodified verifier.

Under the stated global-rank and production-solver assumptions, the modeled
expected operation counts are \LoneModelRange, \LthreeModelRange, and
\LfiveModelRange.  Replacing the field primitive charged by those models with
validated signed AXN circuits gives \LoneAXNRange, \LthreeAXNRange, and
\LfiveAXNRange.  These arithmetic-only figures leave
\AXNNominalHeadroomRange{} bits below NIST's nominal gate references and
\AXNFullAESGapRange{} bits below the released full-key-space AES calibration
circuits in the same basis.

The exact and conditional claims should not be merged.  The quotient, affine
reduction, stable kernel, rejection identities, exact-uniform target sampler,
probability inequalities, and final-verifier wrapper are mathematical or
deterministic public statements.  The headline exponents additionally depend
on global polar spectra that were sampled rather than certified and on
production MQ behavior extrapolated from published models and smaller
experiments.  The circuit conversion prices dominant arithmetic, not a
complete solver.  The strongest unconditional conclusion is therefore that public-matrix
symmetry strictly reduces the quadratic feature space of every symmetric
$q=19$ public key.  For each of the nine parameter shapes, every tested
full-size instance also passes the public preflight and yields the complete
residual system analyzed here.  Turning the reported arithmetic budgets into a
definitive end-to-end security classification requires global rank
certification and production-scale solver measurements.

For a redesign, the immediate lesson is simple: count distinct quadratic
polynomials after every public symmetry, not the labels used to name them.

\appendix
\setcounter{section}{0}
\renewcommand{\thesection}{\Alph{section}}
\renewcommand{\theHsection}{app.\Alph{section}}


\section{Proofs Deferred from the Main Text}\label{app:deferred-proofs}

The main text states the probability and structural results at the point where
an attacker uses them.  This appendix collects the longer proofs in the same
order.

\subsection{Stable-lift kernel}\label{app:proof-stable-kernel}

\begin{proof}[Proof of \Cref{thm:stable-kernel}]
Choose an $\F_q$-basis of $A$.  The displayed spanning family in
\eqref{eq:Uw} then has $m_1d^2$ elements, proving the dimension bound.  For
$\zeta\in A$,
\[
 \bigl(w^\transpose\Sigma_\xi P_i\Sigma_\eta\bigr)\Sigma_\zeta
 =w^\transpose\Sigma_\xi P_i\Sigma_{\eta\zeta},
\]
so $\mathcal U_w$ is closed under right multiplication by $A$.

Consider a cross term produced by substituting
$u_j=\Sigma_{R_j}u+\Sigma_{\Gamma_j}w$ into a feature.  A term with $w$ on the left
already has row coefficient
$w^\transpose\Sigma_\xi P_i\Sigma_\eta$.  A term with $w$ on the right can be
transposed; using $S^\transpose=S$, $P_i^\transpose=P_i$, and commutativity of
$A$, its row coefficient has the same form with the two algebra labels
reversed.  Thus every row of $\mathcal F_{R,\mathbf v}$ lies in $\mathcal U_w$.
Consequently every row of
$C_{R,\mathbf v}=W_R\mathcal E\mathcal F_{R,\mathbf v}$ lies there as well.  Right-$A$
stability then puts every row of the orbit matrix $\mathcal O_{R,\mathbf v}$ in
$\mathcal U_w$.

If $\rank\mathcal O_{R,\mathbf v}=m_1d^2$, then
\[
 m_1d^2=\dim\operatorname{row}(\mathcal O_{R,\mathbf v})
 \le\dim\mathcal U_w\le m_1d^2,
\]
so $\operatorname{row}(\mathcal O_{R,\mathbf v})=\mathcal U_w$.  Since the rows of
$\mathcal F_{R,\mathbf v}$ belong to this row space,
$\mathcal F_{R,\mathbf v}\ker\mathcal O_{R,\mathbf v}=0$.  To see $A$-stability of the
kernel, let $x\in\ker\mathcal O_{R,\mathbf v}$ and $\zeta\in A$.  Every product
$S^a\zeta$ is an $\F_q$-linear combination of the chosen basis powers, so
$C_{R,\mathbf v}\Sigma_{S^a}\Sigma_\zeta x=0$ for all $a$; hence
$\Sigma_\zeta x\in\ker\mathcal O_{R,\mathbf v}$.  Rank--nullity gives
$\dim_{\F_q}H_{R,\mathbf v}=N_0-m_1d^2$.  When $A\cong\F_{q^\ell}$, $H_{R,\mathbf v}$ is an
$A$-vector space, and division by $\ell$ gives
$\dim_A H_{R,\mathbf v}=n-m_1\ell$.
\end{proof}

\subsection{Stable affine slices}\label{app:proof-stable-slices}

\begin{proof}[Proof of \Cref{cor:stable-slices}]
The first two ranks in \eqref{eq:stable-preflight} give the full-lift case of
\Cref{prop:affine-reduction}; in particular, every target determines a
nonempty affine space $a(t)+\ker C_{R,\mathbf v}$.  The orbit matrix contains
$C_{R,\mathbf v}$ as its first block, so
$H_{R,\mathbf v}\subseteq\ker C_{R,\mathbf v}$.  Hence
$a(t)+y+L\subseteq a(t)+\ker C_{R,\mathbf v}$ whenever
$y\in\ker C_{R,\mathbf v}$.  Because
$L\subseteq H_{R,\mathbf v}$, \Cref{thm:stable-kernel} gives
$\mathcal F_{R,\mathbf v}L=0$, and therefore $\mathcal F_{R,\mathbf v}u$ is constant on that
slice.

By \Cref{lem:spectral-quotient}, the unordered power-pair features can be
changed invertibly to eigenblock coordinates.  Because $Q_R$ has full column
rank, row reduction of the complete output system isolates $K$ rows whose
homogeneous quadratic parts are an invertible linear image of
$z_{\rm sym}$.  Applying the inverse image change to those rows expresses the
quadratic parts in eigenblock coordinates.  Full rank of the quadratic feature
map restricted to $L$ ensures that this restriction loses no eigenblock
component.  Translating a diagonal feature introduces only
quadratic, linear, and constant terms from one block; translating an
off-diagonal feature introduces terms from its two blocks.  If $h_a$ is a
homogenizer for block $a$, multiplication by the unique powers of the $h_a$
turns these affine polynomials into exact multidegrees $2\mathbf e_a$ and
$\mathbf e_a+\mathbf e_b$.

Distinct unordered multidegrees have disjoint monomial supports.  Each
component is generated by the $m_1$ base forms and therefore has rank at most
$m_1$.  The total restricted rank is the sum of these component ranks.  Since
there are $\binom{\ell+1}{2}$ components and the total is
$K=m_1\binom{\ell+1}{2}$, every component must have rank exactly $m_1$.
\end{proof}

\subsection{Accepted targets for the \texorpdfstring{$\ell=2$}{l=2} filter}\label{app:proof-l2-filter}

\begin{lemma}[Quadratic fiber bound]\label{lem:quadratic-fiber-bound}
Let $q$ be odd and let $g:\F_q^N\to\F_q^K$ be quadratic.  For each
$b\in\F_q^K\setminus\{0\}$, let $r(b)$ be the polar rank of the scalar
quadratic $b\cdot g$, put $r_*=\min_{b\ne0}r(b)$, and set
$\alpha=1-q^{-K}$.  Then every target $z\in\F_q^K$ satisfies
\begin{equation}\label{eq:quadratic-fiber-bound}
 \left|\,|g^{-1}(z)|-q^{N-K}\right|
 \le \alpha q^{N-r_*/2}.
\end{equation}
More generally, if $W\subseteq\F_q^N$ is an affine subspace of codimension
$c$, then
\begin{equation}\label{eq:quadratic-affine-upper}
 |\{x\in W:g(x)=z\}|
 \le q^{N-c-K}+\alpha q^{N-r_*/2}.
\end{equation}
\end{lemma}

\begin{proof}
Fix a nontrivial additive character $\psi$ of $\F_q$.  Fourier inversion gives
\[
 |g^{-1}(z)|
 =q^{-K}\sum_{b\in\F_q^K}\psi(-b\cdot z)
   \sum_{x\in\F_q^N}\psi(b\cdot g(x)).
\]
The $b=0$ term is $q^{N-K}$.  For $b\ne0$, let $r=r(b)$.  An invertible
linear change of variables splits the polar form into a nondegenerate
$r$-dimensional part and a radical of dimension $N-r$.  If the linear part of
$b\cdot g$ is nonzero on the radical, summing over the radical gives zero.
Otherwise the radical contributes $q^{N-r}$, while diagonalizing the
nondegenerate quadratic part reduces its sum to a product of one-dimensional
Gauss sums of absolute value $q^{1/2}$.  In either case the inner sum has
absolute value at most $q^{N-r/2}$~\cite{Lam2005}.  Summing the $q^K-1$
nonzero Fourier coefficients proves \eqref{eq:quadratic-fiber-bound}.

Write $W=a+U$, where $U$ has codimension $c$.  Translation does not change a
quadratic polynomial's polar form.  Restricting a bilinear form to
$U\times U$ lowers its rank by at most $2c$: after extending a basis of $U$ to
one of $\F_q^N$, deleting the last $c$ rows and columns changes matrix rank by
at most $2c$.  Thus every nonzero scalar combination of $g|_W$ has polar rank
at least $r_*-2c$.  If $r_*\ge2c$, the same Fourier calculation on the
$(N-c)$-dimensional affine space $W$ bounds each nonzero character sum by
\[
 q^{(N-c)-(r_*-2c)/2}=q^{N-r_*/2}.
\]
If $r_*<2c$, the trivial bound $q^{N-c}$ is already at most
$q^{N-r_*/2}$.  This proves \eqref{eq:quadratic-affine-upper} in both cases.
\end{proof}

\begin{proof}[Proof of \Cref{prop:l2-rejection}]
By \Cref{lem:quadratic-fiber-bound}, every target fiber has between
$q^{N-K}-\alpha q^{N-r_*/2}$ and
$q^{N-K}+\alpha q^{N-r_*/2}$ roots.  A rejected root lies on
the intersection of some $t_{\rm rej}$ block-symmetry hyperplanes.  Each such
intersection is an affine subspace of codimension $t_{\rm rej}$, so
\eqref{eq:quadratic-affine-upper} bounds the number of roots it contains by
$q^{N-t_{\rm rej}-K}+\alpha q^{N-r_*/2}$.  Subtracting the union bound
over the $B$ intersections from the full-fiber lower bound gives
\[
 q^{N-K}(1-Bq^{-t_{\rm rej}})
 -\alpha q^{N-r_*/2}(1+B).
\]
This is positive exactly under the sufficient inequality
\eqref{eq:l2-rejection}.

Every fiber has size at most
$q^{N-K}(1+\alpha q^{K-r_*/2})$.  Since
$|\mathcal X_{\rm acc}|=a_n q^N$, double-counting accepted assignments by their images
proves \eqref{eq:l2-accepted-target}.  Finally, let
$T_z=|g^{-1}(z)|$ and let $R_z$ be the number of rejected roots.  The accepted
fraction is $1-R_z/T_z$.  Applying the union-bound upper estimate to $R_z$ and
the character-sum lower estimate to $T_z$ gives
\eqref{eq:l2-accepted-fraction}.
\end{proof}

\subsection{Root collisions and the rank-spectrum bound}\label{app:proof-coset-success}

\begin{proof}[Proof of \Cref{thm:coset-success}]
Average over the $q^{\dim V-h}$ cosets and $q^K$ targets.  Every point of $V$
contributes once to the first moment, which gives \eqref{eq:EZ}.  For an
ordered pair of distinct points in one coset, write the second point as $x+y$
with $y\in L\setminus\{0\}$.  The collision equation is
\[
 g(x+y)=g(x)
 \quad\Longleftrightarrow\quad
 T_y(x)=-(g(y)-g(0)).
\]
It has $\kappa(y)q^{\dim V-r(y)}$ solutions.  Dividing the total number of
ordered collisions by the number of target--coset pairs proves
\eqref{eq:factorial-moment}.

For the rank-spectrum identity, note that
$q^{-r(y)}=q^{-K}|\ker T_y^*|$ and double-count pairs $(b,y)$ satisfying
$b\cdot T_y=0$.  The term $b=0$ contributes $q^h-1$.  For $b\ne0$, the number
of nonzero $y\in L$ annihilated by $b\cdot T_y$ is
$q^{h-e_L(b)}-1$.  Therefore
\[
 \bar\eta
 =q^{-K}\left((q^h-1)+\sum_{b\ne0}(q^{h-e_L(b)}-1)\right)
 =\mu(1+\Theta_L(g))-1.
\]
The inequality $\eta_0\le\bar\eta$ is immediate.

Using $\mathbb E[Z^2]=\mu(1+\eta_0)$ gives
\[
 \operatorname{Var}(Z)
 =\mu(1+\eta_0-\mu)
 \le\mu^2\Theta_L(g),
\]
and the second-moment inequality proves \eqref{eq:coset-success}.  Finally,
$e_*\ge h-\Delta$ implies
$\Theta_L(g)\le(q^K-1)q^{-h+\Delta}<q^{\tau+\Delta}$.
\end{proof}

\subsection{Unique-root trials and the planted distribution}\label{app:proof-singleton-coset}

\begin{proof}[Proof of \Cref{cor:singleton-coset}]
For every nonnegative integer $z$,
\[
 \mathbf 1_{z=1}\ge z-z(z-1),\qquad
 \mathbf 1_{z>0}\ge z-\binom z2,\qquad
 \mathbf 1_{z\ge2}\le\binom z2.
\]
Apply these inequalities to $Z$ and use
\eqref{eq:factorial-moment} together with $\eta_0\le\bar\eta$.
Dividing the upper bound on $\Pr[Z\ge2]$ by
\eqref{eq:bonferroni-success} proves \eqref{eq:conditional-multiple}.

The planted law size-biases a target--coset pair by $Z$.  Hence
\[
 \Pr_{\mathsf P}[Z\ge2]
 =\frac{\mathbb E[Z\mathbf 1_{Z\ge2}]}{\mu}
 \le\frac{\mathbb E[Z(Z-1)]}{\mu}=\eta_0.
\]
Conditioned on $Z=1$, both $\mathsf U$ and $\mathsf P$ are uniform on the same
set of target--coset pairs.  If two distributions have the same conditional
law on an event, their total variation distance is at most the larger mass they
assign to its complement.  Here the complement is $\{Z\ge2\}$.  The same indicator inequality gives the sharper exact-parameter lower bound
$\Pr[Z>0]\ge\mu(1-\eta_0/2)$.  Combining it with
$\Pr[Z\ge2]\le\mathbb E[Z(Z-1)]/2=\mu\eta_0/2$ gives
\[
 \Pr_{\mathsf U}[Z\ge2]
 \le \frac{\eta_0}{2-\eta_0}\le\eta_0
 \qquad(\eta_0<1).
\]
The displayed size-bias bound gives
$\Pr_{\mathsf P}[Z\ge2]\le\eta_0$.  Both are at most $\bar\eta$, proving the
total-variation claim.
\end{proof}

\subsection{Full-space target availability}\label{app:proof-full-space}

\begin{proof}[Proof of \Cref{cor:full-space}]
Take $L=V$ in \Cref{thm:coset-success}.  Then $e_L(b)$ is the polar rank and
$\Theta_L(g)\le(q^K-1)q^{-r_*}=\alpha q^{K-r_*}$.  For the last statement, \eqref{eq:quadratic-fiber-bound} gives
\[
 |g^{-1}(z)|\ge q^{N-K}-\alpha q^{N-r_*/2}>0
\]
when $r_*\ge2K$.
\end{proof}

\section{Dimension Formulas}\label{app:dimensions}

For a Version~2.3 row,
\[
 n=v+o,
 \qquad m_1=\left\lceil\frac{or}{\ell}\right\rceil,
 \qquad M=or\ell,
 \qquad N_0=n\ell,
 \qquad N_{\rm raw}=m_1\ell^2r^2,
 \qquad K=m_1\binom{\ell+1}{2}.
\]
For a full structured lift, $\rank C_{R,\mathbf v}=M-K$ and
\[
 \dim\ker C_{R,\mathbf v}=N_0-M+K.
\]
If $\rank\mathcal O_{R,\mathbf v}=m_1\ell^2$, then
\[
 \dim_{\F_q}H_{R,\mathbf v}=N_0-m_1\ell^2,
 \qquad
 \dim_A H_{R,\mathbf v}=n-m_1\ell.
\]
The selected special-XL slice has $A$-dimension $s$, base-field dimension
$h=\ell s$, and retry exponent $\tau=K-h$.

\begin{table}[tbp]
\centering
\caption{Structured-lift dimensions for the six Version~2.3 $\ell=4$ rows.}
\label{tab:dimensions}
\small
\setlength{\tabcolsep}{5pt}
\begin{tabular}{@{}l r r r r r@{}}
\toprule
Parameters & $M$ & $K$ & $\dim_A H$ & $s$ & $\tau$\\
\midrule
$(28,5,4,4)$ & 80  & 50 & 13 & 10 & 10\\
$(28,4,4,5)$ & 80  & 50 & 12 & 10 & 10\\
$(40,7,4,4)$ & 112 & 70 & 19 & 15 & 10\\
$(38,5,4,5)$ & 100 & 70 & 15 & 15 & 10\\
$(50,9,4,4)$ & 144 & 90 & 23 & 19 & 14\\
$(52,6,4,6)$ & 144 & 90 & 22 & 19 & 14\\
\bottomrule
\end{tabular}
\end{table}

\section{Cost-Search Certificate}\label{app:cost-certificates}

For each $m_1\in\{5,7,9\}$ and every admissible integer $s$, the artifact
expands \eqref{eq:hilbert-series} exactly.  Here admissible means
$1\le s\le\min\{\dim_A H,\lfloor K/4\rfloor\}$.  The $K/4$ bounds are
$12,17,22$.  They bind or tie for both $m_1=5$ rows and both $m_1=9$ rows.
For $m_1=7$, the bounds are $17$ for $(40,7,4,4)$ and $15$ for
$(38,5,4,5)$, the latter coming from $\dim_A H$.  The artifact searches to
the larger group bound and then filters the certificate separately for each
parameter set; both $m_1=7$ rows attain the same optimum at $s=15$.
The monomial count
$M_s(\mathbf d)$ is coordinatewise increasing, so any multidegree that could
match or beat the incumbent lies in a finite rectangular box.  By symmetry it
suffices to enumerate sorted multidegrees.  The search checks 3,690, 11,217,
and 27,449 candidates for $m_1=5,7,9$, respectively.
 
The optima, unique up
to block permutation, are
those in \Cref{tab:xl-optima}.  No zero signed coefficient occurs anywhere
in these bounded searches, so using a strictly negative rather than a
non-positive trigger does not change the result.  The optimal coefficients are
\[
 -14{,}677{,}784,\qquad
 -124{,}983{,}335{,}190,\qquad
 -1{,}164{,}044{,}975{,}088{,}360.
\]
The corresponding Macaulay column counts are
\[
 23{,}417{,}049{,}656,
 \quad 44{,}237{,}557{,}981{,}725{,}696,
 \quad 236{,}094{,}291{,}515{,}544{,}000{,}000.
\]
Substitution of the full-rank retry bound
$\mathsf{Retry}_{K,\tau}$ into \eqref{eq:special-xl-cost} gives published-model exponents
128.81713, 171.68562, and 214.12301.  Replacing $b_4$ by the validated signed AXN circuit
$g_4^{\mathrm{AXN}}$ gives converted arithmetic exponents 132.17187,
175.04035, and 217.47775.  The NAND sensitivity gives 133.04031, 175.90879,
and 218.34618.

\section{Same-Basis Ordered-Label Estimator Baseline}\label{app:ordered-baseline}

The comparison with the official Version~2.3 table mixes cost conventions, so
we also recost a reproducible \emph{estimator baseline} under exactly the same
AXN arithmetic basis as the selected attacks.  This baseline feeds the full
ordered-label equation/variable dimensions from the earlier affine-column
accounting into the pinned generic-MQ optimizer and replaces each charged
base-field operation by $g_1^{\mathrm{AXN}}=441$.  An attacker can indeed feed
the unreduced affine-column equations to a generic solver.  What is
counterfactual is the estimator's treatment of the ordered labels as though
they represented independent generic quadratic features: on every symmetric
key, reversed labels define the same polynomial by
\Cref{thm:symmetric-square}.  Thus the table is a controlled historical
estimator comparison, not evidence for a second quadratic space of the stated
ordered-label rank.  It reports every row rather than only the range quoted in
the introduction.

\begin{table}[H]
\centering
\caption{Historical ordered-label generic-MQ estimator and selected attack,
both recosted in the same AXN arithmetic basis.  ``Estimator input'' gives the
equation/variable dimensions of the unreduced affine-column equations.  The
system is executable, but treating its reversed labels as independent generic
quadratic features is counterfactual on a symmetric key.}
\label{tab:ordered-baseline}
\footnotesize
\setlength{\tabcolsep}{4.1pt}
\renewcommand{\arraystretch}{1.08}
\begin{tabular}{@{}c l c r r r@{}}
\toprule
Level & Parameters & \shortstack{Estimator\\input} & \shortstack{Ordered-label\\AXN} &
\shortstack{Selected\\AXN} & Saving\\
\midrule
I   & $(28,5,4,4)$  & $80/132$  & 206.69 & 132.17 &  74.52\\
I   & $(48,16,2,2)$ & $64/128$  & 171.77 & 133.75 &  38.02\\
I   & $(28,4,4,5)$  & $80/128$  & 206.69 & 132.17 &  74.52\\
\midrule
III & $(40,7,4,4)$  & $112/188$ & 277.52 & 175.04 & 102.48\\
III & $(72,24,2,2)$ & $96/192$  & 242.15 & 187.74 &  54.41\\
III & $(38,5,4,5)$  & $100/172$ & 250.29 & 175.04 &  75.25\\
\midrule
V   & $(50,9,4,4)$  & $144/236$ & 345.49 & 217.48 & 128.01\\
V   & $(96,32,2,2)$ & $128/256$ & 311.08 & 239.50 &  71.58\\
V   & $(52,6,4,6)$  & $144/232$ & 347.13 & 217.48 & 129.66\\
\bottomrule
\end{tabular}
\end{table}

The exact unrounded savings are $38.02124$--$129.65538$ bits.  This comparison
still inherits the generic-MQ and special-XL solving models, but it removes the
field-operation and gate-basis mismatch from the comparison itself.

\section{Fixed-Skew Certificate}\label{app:fixed-skew}

For the official $\ell=2$ matrix, direct arithmetic gives
\[
 9I+10S=\begin{pmatrix}0&1\\1&7\end{pmatrix},
 \qquad e_0^\transpose(9I+10S)=e_1^\transpose.
\]
Together with \eqref{eq:fixed-skew-offset}, this proves
\eqref{eq:fixed-skew-identity} without a rank or genericity assumption.  The
remaining public preflight checks $\rank Q_R=K$ and
$\rank C_{R,\mathbf v}=M-K$.

The resulting residual dimensions are $48/112$, $72/168$, and $96/224$.
The reoptimized generic-MQ convention gives exponents 129.44030, 183.43803,
and 237.84310.  The last optimum uses Hashimoto's $a=1$ boundary and the direct
$\operatorname{MQ}(19,1,1)$ solver described in the main text.  Replacing the coarse factor $b_1$ by the validated signed AXN circuit
$g_1^{\mathrm{AXN}}$ gives converted arithmetic exponents 133.74643,
187.74416, and 242.14922.  By \Cref{cor:full-space}, a global polar-rank
lower bound $r$ adds at most
\[
 \log_2\!\left(1+(1-19^{-K})19^{K-r}\right)
\]
bits for target retries.  Requiring positive AXN arithmetic budget below the nominal NIST references gives
thresholds 46, 68, and 89.  The stronger condition $r\ge2K$, namely 96, 144,
and 192, makes every target solvable.  These are the values in \Cref{tab:l2-costs} and
\Cref{tab:fixed-skew-evidence}.

\section{Artifact and Reproduction}\label{app:artifact}

The release contains the complete paper source, bibliography, and a
self-contained evidence directory with:
\begin{itemize}
\item the public-key generator, direct verifier, and official Level-I
known-answer public key and signature;
\item the quotient, affine-feature, structured-offset, stable-kernel, and
rejection checks;
\item all 121 structured preflight records, all 13 structured-slice records,
and 60 fixed-skew preflights;
\item 102,796 structured cross-polar ranks, 21,992 dual directional-derivative
ranks, 3,888 fixed-skew polar ranks, and 29,424 zero-offset filter ranks,
including the exceptional-direction campaigns;
\item exhaustive $\ell=2$ relation enumeration, both rejection certificates,
accepted-target bounds, and generic-MQ costs;
\item the complete special-XL and two-block cost searches, 128 exact
$\ell=4$ and 48 control Macaulay experiments, and the target-generation bound;
\item 110 exhaustive toy checks of the exact collision and rank-spectrum formulas,
including rank-defective and compatibility-defective cases, plus genericity
diagnostics on six actual $\ell=2$ headline cores;
\item an independent rank implementation checked on 2,036 matrices, a
projective-distinctness audit of all 124,788 structured rank-test directions,
and independent scalar evaluators for the released gate circuits and headline
cost conversions.
\end{itemize}
All finite-field ranks are computed exactly over $\F_{19}$.  The release
records hashes of the scripts, public keys, stable kernels, selected slices,
and polar bases, together with a manifest for every file.

\begingroup
\footnotesize
\printbibliography
\endgroup
\end{document}
